\documentclass[11pt]{article}
\usepackage[a4paper,margin=27mm]{geometry}
\usepackage{lmodern}
\usepackage[T1]{fontenc}
\usepackage{amsmath,amssymb,amsthm,mathtools}
\usepackage{microtype}
\usepackage{tikz,tikz-cd}
\usetikzlibrary{arrows.meta,calc,positioning,decorations.markings,fit}
\usepackage{booktabs,array,enumitem,longtable}
\usepackage{xcolor}
\definecolor{inkblue}{RGB}{30,63,97}
\definecolor{softblue}{RGB}{226,237,246}
\usepackage[colorlinks=true,linkcolor=inkblue,citecolor=inkblue,urlcolor=inkblue]{hyperref}
\hypersetup{pdftitle={Delooping the 3-sphere in homotopy type theory},pdfauthor={Jiatong Yang}}
\setlist{itemsep=0.2em,topsep=0.35em}
\setlength{\parskip}{0.25em}
\linespread{1.05}
\newtheorem{theorem}{Theorem}[section]
\newtheorem{proposition}[theorem]{Proposition}
\newtheorem{lemma}[theorem]{Lemma}
\newtheorem{corollary}[theorem]{Corollary}
\theoremstyle{definition}
\newtheorem{definition}[theorem]{Definition}
\newtheorem{construction}[theorem]{Construction}
\theoremstyle{remark}
\newtheorem{remark}[theorem]{Remark}
\newtheorem{problem}{Problem}
\newenvironment{proofsketch}{\par\noindent\textit{Proof sketch.}\ }{\hfill$\square$\par}
\newcommand{\Z}{\mathbb Z}
\newcommand{\RP}{\mathbb{RP}}
\newcommand{\CP}{\mathbb{CP}}
\newcommand{\HP}{\mathbb{HP}}
\newcommand{\HH}{\widetilde H}
\newcommand{\trunc}[2]{\lVert #2\rVert_{#1}}
\newcommand{\mere}[1]{\lVert #1\rVert_{-1}}
\newcommand{\fib}{\operatorname{fib}}
\newcommand{\cofib}{\operatorname{cofib}}
\newcommand{\Map}{\operatorname{Map}}
\newcommand{\id}{\operatorname{id}}
\newcommand{\im}{\operatorname{im}}
\newcommand{\isContr}{\operatorname{isContr}}
\newcommand{\isEquiv}{\operatorname{isEquiv}}
\newcommand{\Conn}{\operatorname{Conn}}
\newcommand{\Pointed}{\mathcal U_*}
\newcommand{\eqp}{\simeq_*}
\newcommand{\hpo}{\mathbin{\cup}^{h}}
\newcommand{\hpb}{\mathbin{\times}^{h}}
\newcommand{\baseN}{\mathcal N}
\newcommand{\Loc}{\operatorname{Local}}
\newcommand{\Sep}{\operatorname{Sep}}
\newcommand{\Moore}[2]{M(\Z/#1,#2)}
\newcommand{\Sym}{\mathfrak S}
\newcommand{\Alt}{\mathfrak A}
\newcommand{\B}{\mathrm B}
\newcommand{\pbmark}{\arrow[dr,phantom,very near start,"\lrcorner"]}
\newcommand{\pomark}{\arrow[dr,phantom,very near end,"\ulcorner"]}
\title{Delooping the 3-sphere in homotopy type theory}
\author{Jiatong Yang\\[2pt]\small Qiuzhen College, Tsinghua University\\\small\texttt{yangjt22@mails.tsinghua.edu.cn}}
\date{14 October 2026}
\begin{document}
\maketitle
\begin{abstract}
We construct, in univalent type theory with higher inductive types, an explicit pointed connected type $B$ with $\Omega B\eqp S^3$. Finite covers of Seifert types yield two fibre sequences with fibre $S^3$. Cellular homology and Moore-space recognition make their total spaces contractible after inverting $3$ and $10$; fibre preservation and arithmetic fracture then glue the localized bases. The construction and a maximal torus $K(\mathbb Z,2)\to_*B$ with fibre $S^2$ are formalized in Cubical Agda. Classical recognition of deloopings of $S^3$ with a maximal torus identifies the resulting homotopy type with $\mathbb{HP}^\infty$. An internal comparison of the induced multiplication with the Cayley--Dickson $H$-space structure remains open.
\end{abstract}
\noindent\textbf{Keywords:} homotopy type theory, delooping, 3-sphere, localization, Cubical Agda.\\
\textbf{MSC 2020:} 55P60 (primary), 55P35, 03B15.

\section{Introduction and route}\label{sec:intro}

The loop type $\Omega(B,b)=(b=b)$ of a pointed connected type carries path concatenation and all its associativity, unit and inverse coherences. An $\infty$-group structure on a pointed type $G$ consists of a pointed connected type $B$ and a pointed equivalence $\Omega B\eqp G$ \cite[\S3]{highergroups}. Constructing such a delooping supplies the entire coherent structure at once. Buchholtz and Rijke construct an $H$-space structure on $S^3$ and the quaternionic Hopf fibration \cite{cayley}. In their discussion of further work they write that they ``do not know how to define a delooping of $S^3$ in homotopy type theory'', adding that classically this would be $\HP^\infty$ \cite[\S6]{cayley}. We construct such a delooping internally.

\begin{theorem}\label{thm:main}
In univalent type theory with higher inductive types, there is an explicit pointed connected type $B$ and a pointed equivalence $\Omega B\eqp S^3$. More precisely, we construct an element of
\[
\sum_{B:\Pointed}\bigl(\Conn(B)\times(\Omega B\eqp S^3)\bigr).
\]
Moreover, $B$ has a maximal torus: a pointed map $K(\Z,2)\to_*B$ with fibre $S^2$ (Theorem~\ref{thm:torus}). Both statements are formalized in Cubical Agda (Appendix~\ref{app:formal}).
\end{theorem}

\paragraph{Contribution.}
The localization theory of Christensen, Opie, Rijke and Scoccola \cite{cors} supplies separated localization and fibre preservation; Scoccola's arithmetic fracture theorem \cite{scoccola} supplies the final reconstruction. The new geometric input consists of explicit finite covers of Seifert types and their circle-bundle maps to $\RP^2$. Cellular homology and Moore-space recognition give the required local contractions, and a common comparison identifies the two localized loop spaces compatibly on their overlap. The sum in Theorem~\ref{thm:main} is untruncated: it specifies $B$, its connectedness and its pointed loop equivalence. Truncated choices of homology representatives and finite CW presentations supply the local contractibility propositions used in the assembly. The delooping is of the underlying pointed sphere. Agreement of its loop multiplication with the Cayley--Dickson multiplication remains open (Problem~\ref{prob:CD}). The maximal torus identifies the classical homotopy type of $B$ with $\HP^\infty$ by the recognition results discussed in \S\ref{sec:outlook}.

\paragraph{Strategy.}
We construct two fibre sequences with common fibre $S^3$ whose total spaces become contractible after inverting $3$ and $10$, respectively. Fibre preservation then identifies the loops of the localized bases with the corresponding localizations of $S^3$. A third fibre sequence provides a common comparison after inverting $30$, and an arithmetic fracture square glues the bases to obtain $B$.

\paragraph{Origins of the construction.}
This project grew out of an attempt to pursue the further directions suggested by Buchholtz and Rijke's Cayley--Dickson construction \cite{cayley}. I initially asked GPT to help prove that the resulting multiplication on $S^3$ is homotopy associative. Repeated attempts to construct an associator directly led to higher-dimensional coherence problems that we did not succeed in resolving by that approach. In separate preliminary work, an obstruction-theoretic argument developed with GPT subsequently established the mere existence of an associator. This changed the nature of the problem, but did not supply the full coherent structure sought here: even an explicit choice of associator would still leave the higher coherences required for an $A_\infty$-structure to be constructed.

I therefore redirected the investigation towards constructing a delooping of $S^3$, so that the coherent structure would be supplied by a loop type rather than assembled one coherence at a time. GPT suggested the localization and arithmetic fracture strategy developed in this paper, and the construction was subsequently formalized with the assistance of GPT and Claude. The resulting delooping concerns the underlying pointed sphere; comparison with the original Cayley--Dickson multiplication remains open (Problem~\ref{prob:CD}).

\subsection*{Conventions}
We write $\Pointed:=\sum_{X:\mathcal U}X$ for pointed types and suppress underlying-type projections. We use classical truncation degrees: $\tau_nX=\trunc nX$, and $X$ is $n$-connected when $\tau_nX$ is contractible. In particular, $\Conn(B):=\isContr(\tau_0B)$ and $S^3$ is $2$-connected. The symbol $\eqp$ denotes a pointed equivalence. For $f:X\to Y$,
\[
\fib(f,y):=\sum_{x:X}(f(x)=y),
\qquad
\fib(f):=\fib(f,y_0)\quad\text{if $f$ is pointed}.
\]
Homology $\HH_*$ is reduced integral cellular homology. For $m,n\geq1$ we put
\[
\Moore mn:=\cofib\bigl(\Sigma^{n-1}d_m:S^n\to S^n\bigr),
\]
where $d_m:S^1\to S^1$ is the degree-$m$ map; its reduced homology is $\Z/m$ in degree $n$ and zero in other degrees.

For $f:B\to D$ and $g:C\to D$, the homotopy pullback is $B\hpb_DC:=\sum_{b:B}\sum_{c:C}(f(b)=g(c))$. A specified commutative square is \emph{cartesian} when its comparison with this type is an equivalence. The homotopy pushout $B\hpo_AC$ of $B\xleftarrow{f}A\xrightarrow{g}C$ has inclusions $\iota_B,\iota_C$ and paths $\mathsf{glue}(a):\iota_B(fa)=\iota_C(ga)$. Cartesian squares are marked by $\lrcorner$ at the northwest corner and pushout squares by $\ulcorner$ at the southeast corner.

Write $L_m:=L_{d_m}$ for localization \emph{inverting $m$}, with unit $\eta^m_X:X\to L_mX$. A type $Z$ is $m$-local if precomposition with $d_m$ is an equivalence on $S^1\to Z$; equivalently, every loop-power map $p\mapsto p^m$ is an equivalence. A map $g$ is an $L_m$-equivalence if $L_mg$ is an equivalence, or equivalently if precomposition with $g$ is an equivalence on maps into every $m$-local type. A type $A$ is $L_m$-connected if $L_mA\simeq1$. The separated localization is $L'_m:=L_{\Sigma d_m}$, with unit $\eta'^m_X$. For pointed simply connected types, $L_m$ inverts $m$ on homotopy groups \cite[Thm.~4.25]{cors}; for example, $\pi_3(L_3S^3)\cong\Z[1/3]$.

\subsection{The route}
Let $R=\RP^2$, $\baseN=\tau_2R$ and $\eta_2:R\to\baseN$ be the truncation unit. The proof has six steps.

\begin{enumerate}[label=(S\arabic*),leftmargin=3em]
\item \emph{Common fibre} (\S\ref{sec:truncation}). $F:=\fib(\eta_2)\simeq S^3$, by the complex Hopf fibration and a truncation-square theorem.
\item \emph{Finite spans} (\S\ref{sec:seifert}--\ref{sec:truncation}). For $c\in\{O,I\}$ there are spans $R\xleftarrow{q_c}A_c\xrightarrow{p_c}C_c$ with $q_c$ a circle bundle, $p_c$ a finite cover with $N_O=3$, $N_I=10$ sheets, and $\pi_2(A_c)=0$, such that the two naturality squares to the spans of truncations are cartesian. Descent along the pushouts $E_c=R\hpo_{A_c}C_c$, $P_c=\baseN\hpo_{\tau_1A_c}\tau_1C_c$ yields fibre sequences $F\to E_c\to P_c$.
\item \emph{Homology and recognition} (\S\ref{sec:homology}--\ref{sec:construction}). $E_c$ and $P_c$ are simply connected, and $L_3E_O\simeq1$, $L_{10}E_I\simeq1$.
\item \emph{Fibre preservation} (\S\ref{sec:construction}). If $f:E\to Y$ with $Y$ simply connected, then $L_m\fib(f,y)\simeq\fib(L_mf,\eta^m_Yy)$. Consequently $\Omega L_mP_c\simeq L_mF$ whenever $L_mE_c\simeq1$.
\item \emph{Fracture} (\S\ref{sec:construction}). For simply connected $X$ and coprime $m,n$, the unit $X\to L_mX\hpb_{L_{mn}X}L_nX$ is an equivalence. We apply this to the $2$-connected type $F$.
\item \emph{Common comparison} (\S\ref{sec:construction}). Pushouts $E_*=E_O\hpo_RE_I$, $P_*=P_O\hpo_{\baseN}P_I$ give a third fibre sequence $F\to E_*\to P_*$ with $L_{30}E_*\simeq1$, and the localized loop equivalences commute with the induced maps.
\end{enumerate}

\paragraph{Assembly.}
Put $Q_O=L_3P_O$, $Q_I=L_{10}P_I$, $T=L_{30}P_*$, with maps $u_O:Q_O\to T$ and $u_I:Q_I\to T$ induced by $P_O\to P_*\leftarrow P_I$. Define
\begin{equation}\label{eq:B}
B=Q_O\hpb_TQ_I=\sum_{x:Q_O}\sum_{y:Q_I}(u_Ox=u_Iy),
\qquad
\begin{tikzcd}[column sep=large,row sep=large]
B \arrow[r] \arrow[d] \pbmark & Q_O \arrow[d,"u_O"] \\
Q_I \arrow[r,"u_I"'] & T.
\end{tikzcd}
\end{equation}
Looping preserves pointed homotopy pullbacks, and (S4)--(S6) identify the loop cospan with $L_3F\to L_{30}F\leftarrow L_{10}F$. Hence
\begin{equation}\label{eq:loop-chain}
\Omega B\simeq\Omega Q_O\hpb_{\Omega T}\Omega Q_I\simeq L_3F\hpb_{L_{30}F}L_{10}F\simeq F\simeq S^3,
\end{equation}
the third equivalence being (S5). Connectedness of $B$ follows from connectedness of $Q_O,Q_I$ and simple connectedness of $T$. The multiplication on $S^3$ is the transport of loop concatenation along the pointed correction $\varepsilon$ of \eqref{eq:loop-chain} (\S\ref{sec:construction}).

\paragraph{The numbers 3 and 10.}
They are the indices $[\Sym_4:D_8]=3$ and $[\Alt_5:D_6]=10$ (dihedral subgroups of orders $8$ and $6$), which count the sheets of the covers in (S2). Their coprimality is the arithmetic input in (S5).

\section{Seifert types, circle bundles and covers}\label{sec:seifert}

\subsection{Finite Seifert types}

Put $C=K(\Z,1)\simeq S^1$, with basepoint $0$, positive loop $h_C$ and multiplication maps $\mu_s:C\to C$ for $s\in\Z$. The pants graph is the HIT $P_0=C\hpo_1C$, with gluing path $\gamma$ and based loops $a=\operatorname{ap}_{\iota_1}(h_C)$, $b=\gamma\cdot\operatorname{ap}_{\iota_2}(h_C)\cdot\gamma^{-1}$. Let $\partial_1,\partial_2:C\to P_0$ be the inclusions and $\partial_3$ the map with loop $(ab)^{-1}$, specified by the fixed circle equivalence. Throughout, a \emph{marked equivalence of diagrams} retains the specified boundary maps, markings of core circles and bundle fibres, and the homotopies between them.


Set $T^2=C\times C$ and define $[s,n]:T^2\to C$ by $[s,n](u,t)=\mu_s(u)+\mu_n(t)$. The addition is that of the Eilenberg--Mac Lane model of $C$.
\begin{definition}\label{def:seifert}
For positive integers $p,q,r$, the finite Seifert type $M(p,q,r)$ is the pushout
\begin{equation}\label{diag:seifert}
\begin{tikzcd}[column sep=large,row sep=large]
\coprod_{i=1}^3 T^2 \arrow[r,"g"] \arrow[d,"f"'] \pomark & \coprod_{i=1}^3 C \arrow[d] \\
P_0\times C \arrow[r] & M(p,q,r),
\end{tikzcd}
\qquad
\begin{gathered}
f_i=(\partial_i,\id_C),\\
(g_1,g_2,g_3)=([1,p],[1,q],[1-r,r]).
\end{gathered}
\end{equation}
\end{definition}
The right vertices are solid-torus cores, with primitive kernel slopes $(p,-1)$, $(q,-1)$ and $(r,r-1)$. Van Kampen \cite[\S8.7]{hott} gives
\begin{equation}\label{eq:seifert-presentation}
\pi_1M(p,q,r)=\langle x,y,z,h\mid h\text{ central},\quad x^p=y^q=z^r=xyz=h\rangle,
\end{equation}
where $x,y,z$ are the filling cores and $h$ is the second-circle loop. The third boundary has image $z^{1-r}$ and $xy=z^{r-1}$.

\subsection{The two finite covers}

The four types and the finite groups used to construct the covers are
\[
\begin{array}{c@{\qquad}cc@{\qquad}cc@{\qquad}c}
 c&C_c&A_c&G_c&H_c&[G_c:H_c]\\ \hline
 O&M(2,3,4)&M(2,2,4)&\Sym_4&D_8&3\\
 I&M(2,3,5)&M(2,2,3)&\Alt_5&D_6&10.
\end{array}
\]
Here $G_c=\pi_1(C_c)/\langle h\rangle$. The subgroup $D_8$ is the normalizer of a four-cycle subgroup in $\Sym_4$, and $D_6$ is the normalizer of a three-cycle subgroup in $\Alt_5$; subscripts denote group orders. Fix these embeddings and finite coset representatives. The actions of the three boundary generators on $H_c\backslash G_c$ have orbit lengths
\begin{equation}\label{eq:orbit-lengths}
\begin{array}{c@{\qquad}ccc}
 &\partial_1&\partial_2&\partial_3\\ \hline
 O&2,1&3&1,2\\
 I&2,2,2,2,1,1&1,3,3,3&5,5.
\end{array}
\end{equation}

\begin{construction}[Coset covers of the Seifert diagram]\label{con:coset-cover}
For $H\leq G$, let $\mathcal Q_H(b)=\fib(\B H\to\B G,b)$ over $\B G=\mathrm{EM}_1G$, with coset fibre $H\backslash G$ and transport by right multiplication. A classifier $\rho:M(p,q,r)\to\B G$ gives the cover
\begin{equation}\label{eq:coset-total}
\widehat M_H=\sum_{m:M(p,q,r)}\mathcal Q_H(\rho m)\longrightarrow M(p,q,r).
\end{equation}
\end{construction}
We construct $\rho$ on \eqref{diag:seifert}. We use pushout flattening: the total space of a family over a pushout is the pushout of its restrictions \cite{hott,cubical}.

Let $(n_1,n_2,n_3)=(p,q,r)$ and $(s_1,s_2,s_3)=(1,1,1-r)$. Choose triangle generators $g_i\in G$ of exact orders $n_i$, with $g_1g_2g_3=1$. The auxiliary base is
\[
\mathcal O=P_0\hpo_{\coprod_i C}\coprod_i\B(\Z/n_i),
\]
with boundary maps $\partial_i$ and the cyclic-generator maps $C\to\B(\Z/n_i)$. Its map to $\B G$ sends $a,b$ to $g_1,g_2$ and the $i$th cyclic generator to $g_i$; the product relation supplies the third boundary comparison.

For a filling use the central extension
\[
A(n,s)=\Z^2/\langle(n,-s)\rangle
 =\langle v,h\mid [v,h]=1,\ v^n=h^s\rangle.
\]
The fibres of $\B A(n,s)\to\B(\Z/n)$ form an oriented circle family, trivialized on the boundary by the lifts $v,h$. For $\gcd(n,s)=1$, the isomorphism $A(n,s)\cong\Z$, $v\mapsto s$, $h\mapsto n$, gives boundary map $[s,n]$. The core maps to $s^{-1}\bmod n$, which is $1$ in all three fillings. Glue to the trivial family on $P_0$; univalence and flattening identify the total space with \eqref{diag:seifert}. Its composite through $\mathcal O\to\B G$ is $\rho$, with regular fibre sent to $1$ and specified boundary comparisons.

Flattening the pullback of $\mathcal Q_H$ expresses \eqref{eq:coset-total} as a pushout of lifted pieces. Over $\B\langle g_i\rangle$, components are indexed by $H\backslash G/\langle g_i\rangle$, with stabilizer
\[
S_t=\{k\in\langle g_i\rangle\mid tkt^{-1}\in H\},
\]
with its inclusion into $\langle g_i\rangle$ and conjugation map to $H$. An orbit of length $\ell$ has stabilizer of order $d=n_i/\ell$, generated by $g_i^\ell$. The extension restricts to $A(d,s_i)$, by $v'\mapsto v^\ell$, $h'\mapsto h$. Circle-orbit charts therefore give the actual boundary map
\begin{equation}\label{eq:boundary-cover}
T^2\longrightarrow T^2,\qquad (u,t)\longmapsto(\mu_\ell(u),t),
\end{equation}
including its comparison with both adjacent vertices.

The lifted diagrams give the following result.
\begin{theorem}[Short covers]\label{thm:short-cover}
There are specified equivalences
\[
\theta_c:A_c\simeq\sum_{m:C_c}\mathcal Q_{H_c}(\rho_c m).
\]
The maps $p_c=\operatorname{pr}_1\theta_c:A_c\to C_c$ have set fibres of cardinality $N_O=3$ and $N_I=10$.
\end{theorem}
\begin{proof}
\emph{Cancellation and retained boundaries.}
A regular filling has multiplicity $1$, so its base vertex
$\B(\Z/1)$ is contractible. After contracting a spanning tree, a filling
whose boundary contains a graph generator $t$ exactly once has word
$u t^{\epsilon}v$, with $\epsilon\in\{1,-1\}$ and $u,v$ in the remaining graph.
The elementary cancellation removes this loop together with its filling cell.
Indeed, for any target $Z$, after fixing a map from the remaining graph,
the extension data are a loop $\ell$ and a path
$u'\ell^{\epsilon}v'=\mathsf{refl}$, where $u',v'$ are the images of $u,v$.
The function $\ell\mapsto u'\ell^{\epsilon}v'$ is an equivalence of the loop type,
so these data are contractible. The pushout mapping property therefore makes
the inclusion of the remaining graph an equivalence.

Reindexing the circle family along this base equivalence gives the
corresponding total-space equivalence. On the remaining graph its restriction
is the specified trivial family. Write $i$ for this total-space equivalence
and $r$ for its inverse. A retained boundary map $b$ is carried to $r b$,
with comparison $i r b\simeq b$ inherited from the inverse law. Thus the comparison includes the whole
torus map. The changes of fibre coordinates are the shears computed below.
These are the admissibility conditions checked for the primitive cancellations
in the finite tables.

We give the finite reduction explicitly. Permutations below compose as functions, with the right factor applied first. The triangle generators are
\[
\begin{array}{c@{\qquad}ccc}
 &g_1&g_2&g_3=(g_1g_2)^{-1}\\ \hline
 O&(12)&(234)&(1432)\\
 I&(12)(34)&(235)&(15342).
\end{array}
\]
In case $O$, number the sheets $0,1,2$ using representatives $1,g_1,g_2$. The three boundary orbit lists are $(0,1),(2)$; $(0,2,1)$; and $(0),(1,2)$. Contract the tree formed by the $a$- and $b$-edges from sheet $0$. The four remaining graph generators, written as words in the original pants loops, are
\[
\xi_1=a^2,\qquad \xi_2=ab,\qquad
\xi_3=bab^{-1},\qquad \xi_4=b^2a^{-1}.
\]
For each boundary write $(w,k;[s,d])$ for the left torus loops $(wh^k,h)$ and the right row $[s,d]$; its filling relation is $(wh^k)^d=h^s$. A label $A_i,B_i,C_i$ means the orbit starting at sheet $i$ for $\partial_1,\partial_2,\partial_3$, respectively. The initial five fillings are
\[
\begin{array}{c@{\quad}c@{\quad}r@{\quad}c}
 &w&k&[s,d]\\ \hline
 A_0&\xi_1&0&[1,1]\\
 A_2&\xi_3&0&[1,2]\\
 B_0&\xi_4\xi_2&0&[1,1]\\
 C_0&\xi_2^{-1}&0&[-3,4]\\
 C_1&\xi_4^{-1}\xi_3^{-1}\xi_1^{-1}&0&[-3,2].
\end{array}
\]
There are exactly two primitive cancellations: $A_0$ eliminates $\xi_1$ by $\xi_1=h$, and $B_0$ eliminates $\xi_2$ by $\xi_2=\xi_4^{-1}h$. With $u=\xi_3$ and $v=\xi_4$, the surviving data are
\[
\begin{array}{c@{\qquad}ccc}
 &A_2&C_0&C_1\\ \hline
 w&u&v&v^{-1}u^{-1}\\
 k&0&-1&-1\\
 {[s,d]}&[1,2]&[-3,4]&[-3,2].
\end{array}
\]
Choose the free basis $U=v^{-1}u^{-1}$, $V=u$, whose inverse is $u=V$, $v=V^{-1}U^{-1}$, and order the boundaries as $(C_1,A_2,C_0)$. Their words become $U,V,(UV)^{-1}$. Precomposing each torus by $(a,t)\mapsto(a,t-\mu_k(a))$ removes its left winding $k$ and changes the right weight to $s-dk$. Writing $\nu_i$ for the ordered filling multiplicities, we obtain
\[
\begin{aligned}
(\nu_1,\nu_2,\nu_3)&=(2,2,4),
& (k_1,k_2,k_3)&=(-1,0,-1),\\
(s_1,s_2,s_3)&=(-3,1,-3),
& (s_i-\nu_ik_i)_{i=1}^3&=(-1,1,1).
\end{aligned}
\]
Finally change the graph trivialization by winding $1$ on $U$ and $0$ on $V$. The induced boundary shear integers are $(1,0,-1)$, so the weights become $(1,1,-3)$. In the old total-space coordinates the final pants generators are
\[
A=Uh=v^{-1}u^{-1}h,\qquad B=V=u,\qquad
(AB)^{-1}=vh^{-1}.
\]
This gives precisely the rows $[1,2],[1,2],[-3,4]$ of $M(2,2,4)$.

For case $I$, the following table records the corresponding finite data. Sheets are numbered $0,\ldots,9$; the graph letters $\xi_i$ are indexed by the non-tree edges in increasing sheet order, with the $a$-edge before the $b$-edge. The chosen tree consists of the $a$-edges at sheets $0,2,3$ and the $b$-edges at sheets $1,3,4,5,7,9$.
\begin{center}
\small
\begin{tabular}{@{}ll@{}}
\toprule
Data & Case $I$ \\
\midrule
$\partial_1$ orbits & $(0,1),(2,4),(3,5),(6,9),(7),(8)$ \\
$\partial_2$ orbits & $(0),(1,2,3),(4,6,7),(5,8,9)$ \\
$\partial_3$ orbits & $(0,1,5,6,2),(3,4,7,9,8)$ \\
Initial graph rank; fillings & $11;\ 12$ \\
Cancelled boundaries & $A_0,A_2,A_3,B_1,B_4,B_5,A_6,C_0,C_3$ \\
Eliminated generators, in order & $\xi_2,\xi_4,\xi_5,\xi_3,\xi_7,\xi_{10},\xi_6,\xi_1,\xi_8$ \\
Surviving generators & $u=\xi_9,\ v=\xi_{11}$ \\
Surviving boundaries & $A_7:(v^{-1}u^{-1},0;[1,2])$ \\
 & $A_8:(u,0;[1,2]),\quad B_0:(v,1;[1,3])$ \\
Final basis and inverse & $U=v^{-1}u^{-1},\ V=u;\quad u=V,\ v=V^{-1}U^{-1}$ \\
Left windings; weights after removal & $(0,0,1);\quad(1,1,-2)$ \\
Final graph gauge & $(0,0)$ \\
\bottomrule
\end{tabular}
\end{center}
There are exactly nine primitive cancellations. The same basis substitution and the torus shears $(0,0,-1)$ give $[1,2],[1,2],[-2,3]$, the rows of $M(2,2,3)$. The graph ranks $4,11$ also follow from $2N_c$ edges and $N_c$ vertices, and the orbit lists give the filling counts $5,12$.

These are equivalences of the untruncated pushouts. The routine moves are tree contraction, cancellation of an identity leg or a primitive loop--cell pair, graph maps with the displayed inverse substitutions, and torus shears. The finite certificates in Appendix~\ref{app:formal} supply the orbit and tree data, the boundary words, the legality of each cancellation, and the terminal integer calculations; their interpretation retains the specified boundary homotopies and torus squares. In particular, the surviving core circles are unchanged, apart from their reordering: $(C_1,A_2,C_0)$ in case $O$ and $(A_7,A_8,B_0)$ in case $I$. The regular fibre $h$ is retained throughout. Composing the diagram equivalences gives
\[
\theta_c:A_c\simeq\sum_{m:C_c}\mathcal Q_{H_c}(\rho_c m).
\]
The coset families have constant finite cardinality, so their projections give the asserted covers.
\end{proof}

\subsection{The twisted circle bundle}

Write $e=4$ in case $O$ and $e=3$ in case $I$. Staging the first two fillings of $M(2,2,e)$ and cancelling their identity legs gives
\begin{equation}\label{eq:staged-twisted}
W=C_x\hpo_{C_h}C_y\quad(h\mapsto x^2,\ h\mapsto y^2),
\qquad
M(2,2,e)\simeq W\hpo_{T^2}C,
\end{equation}
whose remaining boundary has loops $((xy)^{-1},h)$ and right leg $[1-e,e]$. Let $D_{\mathrm{inv}}:C\to\mathcal U$ have fibre $C$ and transport $t\mapsto-t$ on the positive loop. Its total space $K_{\mathrm{inv}}=\sum_{u:C}D_{\mathrm{inv}}(u)$ is the inversion mapping torus. The Nielsen change
\[
x\mapsto j,\qquad y\mapsto jr^{-1},\qquad
j\mapsto x,\quad r\mapsto y^{-1}x
\]
identifies $W\simeq K_{\mathrm{inv}}$: circle and pushout elimination identify the relations $x^2=y^2$ and $jr=r^{-1}j$, including the relation square and inverse homotopy. Now $h=j^2$, $(xy)^{-1}=rj^{-2}$, and the torus equivalence
\[
\sigma(u,v)=(v,u+v),\qquad \sigma^{-1}(a,t)=(t-a,a)
\]
changes the left boundary to $(j^2,r)$ and the right row to $[e,1]$, since $(1-e)v+e(u+v)=eu+v$. Hence
\begin{equation}\label{eq:twisted-normal-form}
M(2,2,e)\simeq K_{\mathrm{inv}}\hpo_{T^2}C,
\qquad T^2\xrightarrow{(j^2,r)}K_{\mathrm{inv}},\quad T^2\xrightarrow{[e,1]}C.
\end{equation}

Let $T:J\to\mathcal U$ be the orientation two-point family on $J=\RP^1$ and $\widetilde J=\sum_{x:J}T(x)$. Fix circle charts with $\widetilde J\to J$ of degree $2$ and $f_A:\widetilde J\simeq C$. The projective-space construction gives $R=\RP^2=J\hpo_{\widetilde J}1$ \cite{cubical}. For a two-element type $T$ with involution $z\mapsto\bar z$, put
\[
\Z_T=\{v:T\to\Z\mid v(\bar z)=-v(z)\},\qquad
\operatorname{SignC}(T)=K(\Z_T,1).
\]
Evaluation at $z:T$ gives $\epsilon_z:\operatorname{SignC}(T)\simeq C$, and $\epsilon_{\bar z}=-\epsilon_z$. Define $\mathcal F_e:R\to\mathcal U$ by $\operatorname{SignC}(T(x))$ on $J$, by $C$ on the cap, and by the seam equivalence
\begin{equation}\label{eq:sign-clutching}
v\longmapsto\epsilon_z(v)+\mu_e(f_A(x,z)),\qquad (x,z):\widetilde J.
\end{equation}
Univalence defines the family from this clutching; its dependence on $T$ supplies the deck coherences. Flattening gives
\[
M_e:=\sum_{b:R}\mathcal F_e(b)\simeq K_{\mathrm{inv}}\hpo_{T^2}C
\]
with exactly the legs in \eqref{eq:twisted-normal-form}. This defines $q_e:M_e\to R$ and, by the specified equivalence, $q_c:A_c\to R$.

The fibre at the cap is $C$, with generator $r$, and transport around the nontrivial base loop is inversion. The source group and bundle character are
\begin{equation}\label{eq:bundle-character}
\pi_1(M_e)=\langle r,j\mid jrj^{-1}=r^{-1},\ j^2=r^e\rangle,
\qquad r\mapsto0,\quad j\mapsto1\quad\text{in }\Z/2.
\end{equation}
The third filling core is unchanged in \eqref{eq:twisted-normal-form}. If $x,y,z$ are the positive loops of the three source cores, the comparison therefore gives
\[
x\longmapsto j,\qquad y\longmapsto jr^{-1},\qquad
z\longmapsto r,\qquad h\longmapsto j^2=r^e.
\]
Thus the bundle fibre is marked by $z$, and the regular Seifert fibre is $h=r^e$.

\begin{proposition}\label{prop:source-pi2}
For $e=3,4$, the circle bundle $M_e\to R$ has $\pi_2(M_e)=0$.
\end{proposition}
\begin{proof}
The proof has three steps: pull the bundle back to $S^2$; construct its classifying map with nonzero clutching winding; and use injectivity on $\pi_2$ and the fibre sequence to deduce vanishing.

\emph{Pullback to the sphere.} Let
\[
\widetilde M_e=\sum_{s:S^2}\mathcal F_e(qs),
\qquad Q:\widetilde M_e\longrightarrow M_e,
\qquad Q(s,v)=(qs,v),
\]
where $q:S^2\to R$ is the orientation cover. The square formed by $Q$, $q$ and the two bundle projections is cartesian. We first compute $\pi_2(\widetilde M_e)$.

In the two-cap chart for this cover, the base is
$\widetilde J\hpo_{\widetilde J\times\{0,1\}}\{0,1\}\simeq\Sigma C$.
Evaluation in the chosen sheet identifies the circle fibres on the equator and the caps with $C$. The two seam equivalences are translations by
\[
s_0(a)=\mu_e(f_A(a)),\qquad
s_1(a)=\mu_{-e}(f_A(\delta a)),
\]
where $\delta$ is the deck map of $\widetilde J\to J$. The circle chart has a specified homotopy $f_A\delta\simeq f_A$, including its vertex paths. It identifies these seams with translations by $\mu_e$ and $\mu_{-e}$ in the common equatorial coordinate (Figure~\ref{fig:two-cap}). In the chosen suspension coordinate the clutching winding is $-2e$; only its nonvanishing is used.

\begin{figure}[htbp]
\centering
\begin{tikzcd}[column sep=huge,row sep=large]
C_0 \arrow[d] & C_{\mathrm{eq}}\times C
  \arrow[l,"{(a,t)\mapsto t+\mu_e(a)}"']
  \arrow[r,"{(a,t)\mapsto t+\mu_{-e}(a)}"]
  \arrow[d,"\operatorname{pr}_1"] & C_1 \arrow[d] \\
1 & C_{\mathrm{eq}} \arrow[l] \arrow[r] & 1
\end{tikzcd}
\caption{The pulled-back circle bundle in two-cap coordinates. The lower homotopy pushout is $\Sigma C\simeq S^2$ and the upper one is $\widetilde M_e$. The seam translations have windings $+e,-e$; cancelling the first cap leaves their difference $(-e)-e=-2e$.}
\label{fig:two-cap}
\end{figure}

\emph{Classifying map and winding.} Put $K=K(\Z,2)$ and let $\omega:C\simeq\Omega(K,0)$ be the pointed circle equivalence, with its specified basepoint comparison. For the path family $U(k)=(0=k)$, transport along a loop $\omega(s)$ is right concatenation by $\omega(s)$. The compatibility of $\omega$ with addition therefore identifies translation by $s$ on $C$ with this transport. Assigning $0$ to the equator and both caps, and $\omega(s_b(a))$ to the $b$th seam, defines a map to $K$ whose pulled-back path family is the circle family above. On cancelling the first cap, the remaining seam is
\[
\omega(\mu_e(a))^{-1}\cdot\omega(\mu_{-e}(a))
 =\omega(\mu_{-2e}(a)),
\]
where the equality is the homotopy supplied by the addition and inverse laws. Thus, with $d=-2e$, the resulting map is
\[
c:\Sigma C\longrightarrow K,
\qquad c(\mathsf{north})=c(\mathsf{south})=0,
\qquad \operatorname{ap}_c(\mathsf{merid}(a))=\omega(\mu_d(a)).
\]
The marked cancellation retains the seam homotopies, so it gives an equivalence of total spaces
\[
\widetilde M_e\simeq\sum_{a:\Sigma C}(0=c(a))
 \simeq\fib(c,0)=:V_e.
\]
The second equivalence reverses the path in each fibre.

\emph{Injectivity and the fibre sequence.} To establish injectivity of $c_*$ on $\pi_2$, let
$\sigma:C\to\Omega(\Sigma C,\mathsf{north})$ be the suspension unit,
$\sigma(a)=\mathsf{merid}(a)\cdot\mathsf{merid}(0)^{-1}$,
with its basepoint comparison, and put
$w=\sigma_*[h_C]\in\pi_1\Omega\Sigma C=\pi_2(\Sigma C)$.
The suspension--loop adjunction gives the pointed homotopy
$\Omega c\circ\sigma\simeq\omega\circ\mu_d$.
Choose the integer coordinate on $\pi_2K$ induced by $\omega$ and the positive generator of $\pi_1C$. Since $(\mu_d)_*[h_C]=d[h_C]$, this homotopy gives
\[
c_*(w)=d=-2e\neq0
\]
in that target coordinate. Also $\Sigma C\simeq S^2$, so $\pi_2(\Sigma C)\cong\Z$, and $\pi_2K\cong\Z$. Under any group isomorphisms with $\Z$, the homomorphism $c_*$ has the form $n\mapsto kn$. Its nonzero value on $w$ forces $k\neq0$, and cancellation in $\Z$ then makes it injective. Only a nonzero image is used; the sign and orientation of this suspension coordinate are not compared with those of the standard sphere.

The fibre sequence for $c$ gives
\[
0=\pi_3K\longrightarrow\pi_2(V_e)
 \longrightarrow\pi_2(\Sigma C)
 \xrightarrow{c_*}\pi_2K.
\]
Exactness and injectivity of $c_*$ give $\pi_2(V_e)=0$. The fibre $V_e$ is connected because $\Sigma C$ is connected and $\pi_1K=0$, so vanishing transfers to $\widetilde M_e$ at every basepoint. Finally $\fib(Q,(b,v))\simeq\fib(q,b)$ is a two-element set; hence $Q$ induces an isomorphism on $\pi_2$. Thus $\pi_2(M_e)=0$ for $e=3,4$, with basepoint transport used only in the proposition expressing vanishing.
\end{proof}

\subsection{Truncation squares and the common fibre}\label{sec:truncation}

\begin{theorem}[Truncation squares]\label{thm:trunc-square}
Let $n\geq0$ and $k=n+1$. If every fibre of $f:X\to Y$ is an $n$-type, its naturality square with $\eta_k$ is cartesian. Explicitly, for the canonical path $c_x:\tau_kf(\eta_kx)=\eta_k(fx)$, the comparison
\begin{equation}\label{eq:trunc-comparison}
X\longrightarrow\sum_{y:Y}\sum_{a:\tau_kX}(\tau_kf(a)=\eta_ky),
\qquad x\longmapsto(fx,\eta_kx,c_x)
\end{equation}
is an equivalence.
\end{theorem}
\begin{proof}
The universe $\mathcal U^{\leq n}$ is an $(n+1)$-type \cite[\S7.1]{hott}, so the classifying map $y\mapsto\fib(f,y)$ of $f$ extends along $\eta_k:Y\to\tau_kY$ to a family $D$ over $\tau_kY$ with $D(\eta_ky)\simeq\fib(f,y)$. The total space $\sum_tD(t)$ is a $k$-type, and truncation induction on $\tau_kX$ and on $\tau_kY$ gives mutually inverse maps $\tau_kX\rightleftarrows\sum_tD(t)$ over $\tau_kY$. The pullback of $\sum_tD(t)\to\tau_kY$ along $\eta_k$ is $\sum_y\fib(f,y)\simeq X$, which gives \eqref{eq:trunc-comparison}.
\end{proof}
Taking fibres in the other direction gives
\begin{equation}\label{eq:unit-fibres}
\fib(\eta_k:X\to\tau_kX,a)\simeq
\fib(\eta_k:Y\to\tau_kY,\tau_kf(a)).
\end{equation}

\begin{lemma}[Pointed Hopf reflection]\label{lem:hopf-reflection}
The pointed Hopf classifying map $h:S^2\to_*K(\Z,2)$ induces a pointed equivalence
\[
\bar h:\tau_2S^2\eqp K(\Z,2),\qquad \bar h\circ\eta_2\sim_*h.
\]
There is an unpointed equivalence
$\fib(\eta_2:S^2\to\tau_2S^2,\eta_2(*))\simeq S^3$.
\end{lemma}
\begin{proof}
The Hopf construction gives $h$ with basepoint fibre $S^3$ \cite[\S8.5]{hott}. Since $K(\Z,2)$ is connected, the propositional assertion that its fibres are $2$-connected extends from the basepoint to every point. Thus $h$ is $2$-connected with a $2$-truncated target, so it exhibits the $2$-truncation of $S^2$. The universal property gives $\bar h$ and its unit comparison; the pointing of $h$ supplies the pointing of $\bar h$ and compatibility of the comparison. Taking the basepoint fibre in this square gives the stated sphere equivalence.
\end{proof}
Applying \eqref{eq:unit-fibres} to the orientation cover $q:S^2\to R$ at level $2$ now gives:
\begin{proposition}[Common fibre]\label{prop:common-F}
For $r_0=q(*)$ and $\baseN=\tau_2R$, there is an equivalence
$F=\fib(\eta_2:R\to\baseN,\eta_2r_0)\simeq S^3$.
\end{proposition}
The base $\baseN$ has $\pi_1=\Z/2$ and $\pi_2=\Z$ with sign action.

For $q_c$ use Theorem~\ref{thm:trunc-square} at level $2$ and the equivalence $\tau_2A_c\simeq\tau_1A_c$ supplied by Proposition~\ref{prop:source-pi2}. For $p_c$ use level $1$. This gives the two mixed cartesian naturality squares on the left of \eqref{diag:mixed-squares}; $\bar q_c$ denotes the induced map to $\baseN$.

\begin{proposition}[Cartesian pushout descent]\label{prop:pushout-descent}
A map of spans whose two naturality squares are cartesian induces cartesian squares from each vertex to the corresponding pushouts.
\end{proposition}
The cartesian squares identify the source span with the total spaces of compatible families over the target span. Glue these families by pushout induction; the flattening lemma identifies their total space with the source pushout, and its restrictions to the vertices give the asserted cartesian squares \cite[\S6.12]{hott}.

Apply descent to $R\xleftarrow{q_c}A_c\xrightarrow{p_c}C_c$ and $\baseN\xleftarrow{\bar q_c}\tau_1A_c\xrightarrow{\tau_1p_c}\tau_1C_c$. Their pushouts are $E_c$ and $P_c$, as in (S2), and the third square below is cartesian:
\begin{equation}\label{diag:mixed-squares}
\begin{tikzcd}[column sep=1.4em,row sep=large]
A_c \arrow[r,"q_c"] \arrow[d,"\eta_1"'] \pbmark & R \arrow[d,"\eta_2"]
&[-.3em] A_c \arrow[r,"p_c"] \arrow[d,"\eta_1"'] \pbmark & C_c \arrow[d,"\eta_1"]
&[-.3em] R \arrow[r] \arrow[d,"\eta_2"'] \pbmark & E_c \arrow[d,"t_c"] \\
\tau_1A_c \arrow[r,"\bar q_c"'] & \baseN
& \tau_1A_c \arrow[r,"\tau_1p_c"'] & \tau_1C_c
& \baseN \arrow[r] & P_c.
\end{tikzcd}
\end{equation}
It identifies the specified fibre of $t_c$ with $F$ and gives fibre sequences $F\to E_c\xrightarrow{t_c}P_c$. This proves (S1) and (S2).

\begin{lemma}[Dihedral-character rigidity]\label{lem:dihedral-character}
Let $Q$ be a group generated by $r,j$ with $jrj^{-1}=r^{-1}$, and let $\chi:Q\to\Z/2$ be a homomorphism satisfying $\chi(r)=0$ and $\chi(j)=1$. For $m\geq3$, let $\epsilon:D_{2m}\to\Z/2$ be the rotation/reflection character. If a homomorphism $\lambda:Q\to D_{2m}$ has nonabelian image, then $\chi=\epsilon\circ\lambda$.
\end{lemma}
\begin{proof}
Put $a=\lambda(r)$ and $b=\lambda(j)$. They do not commute, since they generate the image. The relation $bab^{-1}=a^{-1}$ implies that $a^2=1$ would force commutation. Every reflection squares to $1$, so $a$ is a rotation. Since rotations commute, $b$ is a reflection. The two characters agree on the generators.
\end{proof}

\begin{proposition}\label{prop:seifert-simple}
The types $E_c$ and $P_c$ are simply connected.
\end{proposition}
\begin{proof}
The types $R$, $A_c$ and $C_c$ are connected. For $R$ this follows from its projective-space pushout; for the Seifert types, the connected vertex $P_0\times C$ is joined to each connected filling core by a nonempty boundary torus. Choose a basepoint $a_0:A_c$, use its images as basepoints in $R$ and $C_c$, and use the pushout seam at $a_0$ to compare their images in $E_c$. Van Kampen \cite[\S8.7]{hott} identifies $\pi_1(E_c)$ with the pushout of groups
\[
\pi_1(R)\xleftarrow{\chi_c}\pi_1(A_c)
 \xrightarrow{(p_c)_*}\pi_1(C_c),
\qquad \chi_c=(q_c)_*:\pi_1(A_c)\longrightarrow\Z/2.
\]
The bundle character is surjective, since \eqref{eq:bundle-character} gives $\chi_c(j)=1$. Equivalently, its surjectivity follows from the connected circle fibres of $q_c$ and the low-degree homotopy exact sequence.

Write $L_c=\pi_1(C_c)$ and let
\[
K_c=\bigl\langle\!\bigl\langle
 (p_c)_*(\ker\chi_c)
\bigr\rangle\!\bigr\rangle_{L_c}
\]
be the normal closure in $L_c$. Since $\chi_c$ is surjective, a compatible pair of homomorphisms from the span to a group $L$ is determined by a homomorphism $v:L_c\to L$ for which $v\circ(p_c)_*$ kills $\ker\chi_c$. Equivalently, $v$ kills $K_c$. The two universal properties therefore give
\[
\pi_1(E_c)\cong L_c/K_c.
\]

We identify elements of $K_c$ using the actual core lifts. The coset-cover model has a map to $\B H_c$, inducing
$\lambda_c:\pi_1(A_c)\to H_c$ after transport along $\theta_c$. The marked classifier sends the third target core in case $O$ to the chosen four-cycle rotation in $H_O=D_8$, and the second target core in case $I$ to the chosen three-cycle rotation in $H_I=D_6$. Their based lifts give elements $u_c\in\pi_1(A_c)$ satisfying
\[
(p_O)_*(u_O)=z,\qquad (p_I)_*(u_I)=y,
\qquad \lambda_c(u_c)=\rho_c^{\mathrm{rot}},
\]
where $\rho_c^{\mathrm{rot}}$ denotes the chosen rotation in $H_c$. Indeed, in the pullback of $\B H_c\to\B G_c$, a target loop whose classifier is the image of a specified element of $H_c$ lifts using that element and the classifier comparison. The same construction lifts the marked reflection word, so the image of $\lambda_c$ contains the chosen rotation and reflection. These two elements do not commute in either $D_8$ or $D_6$.

Let $\epsilon_c:H_c\to\Z/2$ be the dihedral character, equal to $0$ on rotations and $1$ on reflections. We claim
\[
\chi_c=\epsilon_c\circ\lambda_c.
\]
Apply Lemma~\ref{lem:dihedral-character} to the generators supplied by the marked bundle comparison \eqref{eq:bundle-character} and the noncommuting rotation and reflection in the image of $\lambda_c$; their mere existence suffices because equality of homomorphisms is a proposition.
In particular,
$\chi_c(u_c)=\epsilon_c(\rho_c^{\mathrm{rot}})=0$.
Hence $z\in K_O$ and $y\in K_I$.

It remains to compute the quotients using \eqref{eq:seifert-presentation}; denote the images of its generators by the same letters. In both cases $x^2=xyz$ implies $x=yz$, and therefore $z=y^{-1}x$.
For $O$, the relation $z=1$ gives $x=y$. Then
\[
y^2=x^2=y^3
\]
forces $y=1$ by cancellation, whence $x=1$ and $h=x^2=1$. Thus all four generators of $L_O/K_O$ are trivial.
For $I$, the relation $y=1$ gives $h=y^3=1$ and $x^2=h=1$. Also $z=y^{-1}x=x$, so $z^5=h$ gives $x^5=1$. Since
\[
x^5=(x^2)^2x=x,
\]
we obtain $x=z=1$. Again all generators are trivial. Consequently $\pi_1(E_c)=1$ in both cases.

The pushout $E_c$ is connected because its vertices are connected and its overlap is inhabited. The vertices $\baseN=\tau_2R$ and $\tau_1C_c$ of $P_c$ are also connected, and $\tau_1A_c$ is inhabited, so $P_c$ is connected. Finally the fibre sequence $F\to E_c\to P_c$ gives the exact segment
\[
\pi_1(E_c)\longrightarrow\pi_1(P_c)\longrightarrow\pi_0(F).
\]
Here $\pi_1(E_c)=1$ and $\pi_0(F)=1$, since $F\simeq S^3$ is connected. Exactness therefore gives $\pi_1(P_c)=1$. Connectedness and triviality of these fundamental groups prove that $E_c$ and $P_c$ are simply connected.
\end{proof}

\section{Moore recognition and homology of the total spaces}\label{sec:homology}

The profiles needed below are $\HH_3(E_O)\cong\Z/3$ and $\HH_2(E_I)\cong\Z/2$, $\HH_3(E_I)\cong\Z/10$, with all other reduced homology groups zero. Moore recognition and localization will turn these profiles into contractions of $L_3E_O$ and $L_{10}E_I$; the degree-two class of $E_I$ requires a first Moore attachment to produce a $2$-connected cofibre. The dimension bound in Moore recognition converts finite connectivity estimates into inverse maps by bounded cellular lifting.

\subsection{Finite Moore recognition}

\begin{definition}
A \emph{finite CW presentation of dimension at most $d$} consists of finite sets $I_q$ for $0\leq q\leq d$, skeleta $X_{-1}=\varnothing,X_0,\ldots,X_d$, attaching maps $\alpha_q:I_q\times S^{q-1}\to X_{q-1}$, and specified skeletal pushout squares
\[
\begin{tikzcd}[column sep=large,row sep=large]
I_q\times S^{q-1} \arrow[r,"\alpha_q"] \arrow[d,"\operatorname{pr}_1"'] \pomark & X_{q-1} \arrow[d] \\
I_q \arrow[r] & X_q.
\end{tikzcd}
\]
Here $S^{-1}=\varnothing$, so $X_0\simeq I_0$; the lower-left vertices are the centres of the attached cones. Write $\mathrm{CW}_{\leq d}$ for the type of these presentations and $|C|=X_d$ for the realization of $C$. The assertion that $X$ \emph{merely admits a finite CW model of dimension at most $d$} is
\[
\mere{\sum_{C:\mathrm{CW}_{\leq d}}(|C|\simeq X)}.
\]
A finite HIT presentation specifies finitely many constructors; a finite CW presentation additionally organizes the constructors into these sphere-attaching skeleta. The finite CW bounds used below include this skeletal data.
\end{definition}

With the Moore-space convention of the introduction, write $\delta_{m,n}:=\Sigma^{n-1}d_m:S^n\to S^n$ for $m,n\geq1$. Naturality of the cellular suspension isomorphisms and $\HH_1(d_m)=\times m$ give $\HH_n(\delta_{m,n})=\times m$ in the standard sphere coordinates.

A map is $k$-connected when all its fibres are $k$-connected. For a pointed $(n-1)$-connected CW type $X$ and $n\geq1$, the cellular Hurewicz theorem gives a natural isomorphism
\[
\pi_n(X,x)^{\mathrm{ab}}\cong\HH_n(X).
\]
Abelianization is redundant for $n\geq2$. The cellular Hurewicz map, its naturality, and this isomorphism are supplied directly by the pinned Cubical library (Appendix~\ref{app:formal}). For the prespectrum-homology formulation in HoTT, see \cite[Theorem~3.12]{hurewiczhott}.

\begin{theorem}[Four-dimensional Moore recognition]\label{thm:moore-recognition}
Let $m>0$ and let $X$ be a $2$-connected type merely admitting a finite CW model of dimension at most $4$. If $\HH_3(X)\cong\Z/m$ and $\HH_4(X)=0$, then
\[
\mere{X\simeq\Moore m3}.
\]
\end{theorem}
\begin{proof}
Write $M=\Moore m3$ and fix the given group isomorphism $\theta_X:\HH_3(X)\cong\Z/m$. Choose a finite CW model of $X$ inside the propositional truncation of the conclusion; all subsequent choices remain inside that truncation.

\emph{Step 1: the calibrated representative.}
Choose a point $x:X$, merely available by connectedness, and put $h=\theta_X^{-1}(\overline1)$. Hurewicz gives a group isomorphism
\[
\mathfrak h_X:\pi_3(X,x)\xrightarrow{\cong}\HH_3(X).
\]
Since $\pi_3(X,x)$ is the set truncation of pointed sphere maps, there merely exists $s:S^3\to_*(X,x)$ with $\mathfrak h_X([s])=h$. Let $u_3\in\HH_3(S^3)$ be the standard generator. Hurewicz naturality and the suspension degree calculation give
\[
\mathfrak h_X([s\circ\delta_{m,3}])
=\HH_3(s\circ\delta_{m,3})(u_3)
=m\HH_3(s)(u_3)=mh=0.
\]
Injectivity of $\mathfrak h_X$ gives $[s\circ\delta_{m,3}]=0$. Effectiveness of set truncation turns this equality into a merely existing pointed nullhomotopy. Forgetting its pointing gives $v:s\circ\delta_{m,3}\sim\operatorname{const}_x$. Cofibre recursion now defines $f:M\to X$: it sends the cone point to $x$, restricts to $s$ on the bottom inclusion $i:S^3\to M$, and sends the cofibre path at $z$ to $v_z^{-1}$. In particular, it retains a specified homotopy $f i\sim s$.

For completeness, the homology coordinate on $M$ is calibrated on this bottom inclusion. The degree calculation and the cofibre sequence identify
\[
\HH_3(M)\cong\operatorname{coker}(\Z\xrightarrow{\times m}\Z).
\]
The quotient map induced by $a\mapsto\HH_3(i)(a u_3)$ is an isomorphism $\Z/m\to\HH_3(M)$. Its inverse $\theta_M$ satisfies $\theta_M(\HH_3(i)(u_3))=\overline1$. Naturality and the restriction homotopy give
\[
\HH_3(f)(\HH_3(i)(u_3))=\HH_3(s)(u_3)=h.
\]
Thus $\theta_X\HH_3(f)\theta_M^{-1}$ fixes $\overline1$. A homomorphism out of $\Z/m$ is determined by this element, so this composite is the identity and $\HH_3(f)$ is an isomorphism. This argument includes $m=1$.

\emph{Step 2: connectivity of the actual cofibre.}
Both $M$ and $X$ are $2$-connected. Hurewicz naturality identifies $\pi_3(f)$ with $\HH_3(f)$ at every point of $M$. The induced map $\tau_3f:\tau_3M\to\tau_3X$ is therefore an equivalence by Whitehead for truncated types \cite[\S8.8]{hott}: the component types are contractible, $\pi_1$ and $\pi_2$ vanish, $\pi_3(f)$ is an isomorphism, and all higher groups of these $3$-types vanish. The truncation units are $3$-connected. Applying connectivity cancellation to
\[
\eta_3^X f=(\tau_3 f)\eta_3^M
\]
yields the $2$-connectedness of $f$ used below. Pushout connectivity then makes $K=\cofib(f)$ $3$-connected.

The finite dimension bound gives $\HH_5(X)=0$. The cellular differential of $M$ in degrees $4$ and $3$ is $\times m:\Z\to\Z$, so $m>0$ gives $\HH_4(M)=0$. Choose the cellular approximations for the actual map $f$ under truncation. Its cofibre homology sequence contains
\[
\begin{gathered}
\HH_4(X)\longrightarrow\HH_4(K)\longrightarrow\HH_3(M)
 \xrightarrow{\HH_3(f)}\HH_3(X),\\
\HH_5(X)\longrightarrow\HH_5(K)\longrightarrow\HH_4(M).
\end{gathered}
\]
Exactness and $\HH_4(X)=0$ give $\HH_4(K)=\HH_5(K)=0$. Hurewicz for the $3$-connected CW type $K$ identifies $\pi_4(K)$ with $\HH_4(K)$; its vanishing makes $\tau_4K$ contractible by truncated Whitehead, so $K$ is $4$-connected. Hurewicz in degree $5$ similarly gives $\pi_5(K)=0$ and makes $K$ $5$-connected. These are two finite connectivity improvements.

\emph{Step 3: finite cellular lifting.}
Let $k_0:K$ be the cone point and $\iota:X\to K$ the other cofibre inclusion. In the cofibre square the maps $M\to1$ and $f:M\to X$ are both $2$-connected. Blakers--Massey \cite{bm} gives a $4$-connected comparison
\[
M\longrightarrow P:=\sum_{x:X}(k_0=\iota x),
\qquad a\longmapsto(fa,\mathsf{glue}(a)).
\]
Since $K$ is $5$-connected, every path type $k_0=\iota x$ is $4$-connected. Thus $P\to X$ is $4$-connected, and composition shows that every fibre of $f$ is $4$-connected.

We use the finite cellular lifting lemma in its precise form: a family of $(d-1)$-connected types over a type merely admitting a finite CW model of dimension at most $d$ merely has a section. To prove it, induct over the finite skeleta. Over a $q$-cell the boundary section, after transport to the fibre at the cell centre, is a map from $S^{q-1}$ to that fibre. Connectedness gives a merely existing cone on this map for $q\leq d$; for $q=0$ it gives a point. Finite choice combines the fillers for the finitely many cells at each stage, and dependent pushout elimination extends the section. All stages and the finite presentation are eliminated into propositional truncation.

Apply this lemma with $d=4$ to $x\mapsto\fib(f,x)$ over $X$, using its $3$-connectedness. It gives, merely, a section $\sigma(x)=(g(x),\epsilon_x)$, where $\epsilon_x:f(g(x))=x$. For $a:M$, the type
\[
\sigma(fa)=(a,\mathsf{refl})\quad\text{in }\fib(f,fa)
\]
is $3$-connected, since $\fib(f,fa)$ is $4$-connected. The same lemma over the finite $4$-dimensional type $M$ gives, merely, a dependent family of these paths. Projection to $M$ gives $g f\sim\id_M$, while $\epsilon$ gives $f g\sim\id_X$. Equivalently, the second lifting uses relative cells of dimension at most $5$ in a cellular cylinder on $M$. Hence $f$ is an equivalence. For this fixed $f$, both truncated lifting choices eliminate into the proposition $\isEquiv(f)$; inverting $f$ and retaining the original truncation gives $\mere{X\simeq M}$.
\end{proof}

\begin{remark}
The cofibre $K$ merely has a finite CW model of dimension at most $5$: coning the $4$-cells of $M$ can introduce $5$-cells. The proof uses $\HH_4(K)=\HH_5(K)=0$ and the dimension bounds on $M$ and $X$.
\end{remark}

For the localizations $L_m$ fixed in the introduction, Moore spaces have the following contractions.

\begin{corollary}[Local Moore contractions]\label{cor:moore-null}
For $n\geq1$ and positive $d\mid m$, the type $\Moore dn$ is $L_m$-connected. In particular,
\begin{equation}\label{eq:needed-moore-null}
L_3\Moore33\simeq1,\qquad
L_{10}\Moore22\simeq1,\qquad
L_{10}\Moore{10}3\simeq1.
\end{equation}
\end{corollary}
\begin{proof}
Write $m=dk$ and let $P_j$ denote the $j$th power map on a loop type. The identities
\[
P_m=P_dP_k=P_kP_d
\]
show that $P_d$ and $P_k$ commute; if $P_m$ is invertible, so are both factors, with inverses $P_kP_m^{-1}$ and $P_dP_m^{-1}$ respectively. The loop-power criterion \cite[Lemma~4.2]{cors} therefore makes every $m$-local type $d$-local. Apply the same argument to iterated loop types, using closure of local types under identity types. For every $m$-local $Z$ and $z:Z$, precomposition by $\delta_{d,n}=\Sigma^{n-1}d_d$ is an equivalence on $\Map_*(S^n,(Z,z))$. The pointed cofibre mapping property identifies $\Map_*(\Moore dn,(Z,z))$ with its fibre over the constant map, hence makes it contractible. This is the implication (1)$\Rightarrow$(2) of \cite[Theorem~3.6]{cors} applied to $\delta_{d,n}$. Evaluation $\Map(\Moore dn,Z)\to Z$ is consequently an equivalence for every such $Z$. The local mapping criterion \cite[Lemma~2.9]{cors} gives $L_m\Moore dn\simeq1$.
\end{proof}

\begin{lemma}[Two-stage local connectedness]\label{lem:two-stage-null}
For $a:A\to E$, if $A$ and $\cofib(a)$ are $L_m$-connected, then $E$ is $L_m$-connected.
\end{lemma}
\begin{proof}
For an $m$-local type $Z$, the cofibre mapping property gives
\[
\Map(\cofib(a),Z)\simeq
\Map(E,Z)\hpb_{\Map(A,Z)}Z.
\]
The constant-map inclusion $Z\to\Map(A,Z)$ is an equivalence because $A$ is $L_m$-connected. Thus restriction $\Map(\cofib(a),Z)\to\Map(E,Z)$ is an equivalence. The composite with $Z\to\Map(\cofib(a),Z)$ is also an equivalence and is the constant-map inclusion for $E$. Apply \cite[Lemma~2.9]{cors}.
\end{proof}

\subsection{Cellular homology and naturality}

For a finite CW presentation, the ordinary cellular chains $C_q(X)=\Z[I_q]$ for $q\geq0$ are free on the $q$-cells, with differential given by attaching degrees and $C_{-1}(X)=0$. The augmentation $\epsilon_X:C_0(X)\to\Z$ is specified separately. We use independence of the presentation, functoriality for arbitrary maps between the realized types, homotopy invariance, vanishing above the CW dimension, and the integral Hurewicz isomorphism stated above. Below, $H_q=\HH_q$ for $q>0$.

Sweeping the cells of the overlap along the gluing interval gives the ordinary double mapping cylinder complex for a cellular span $B\xleftarrow f A\xrightarrow g C$:
\begin{equation}\label{eq:DMC}
\begin{aligned}
\mathcal D_q&=C_q(B)\oplus C_q(C)\oplus C_{q-1}(A) &&(q\geq0),\\
\partial(b,c,a)&=(\partial b+f_*a,\partial c-g_*a,-\partial a) &&(q\geq1).
\end{aligned}
\end{equation}
Thus the low-degree differential and the single augmentation are
\[
\begin{aligned}
\mathcal D_0&=C_0(B)\oplus C_0(C),\\
\partial_1(b,c,a)&=(\partial b+f_*a,\partial c-g_*a),\\
\epsilon_{\mathcal D}(b,c)&=\epsilon_B(b)+\epsilon_C(c).
\end{aligned}
\]
The vertex inclusions and the projection $(b,c,a)\mapsto a$ give the reduced Mayer--Vietoris sequence
\begin{equation}\label{eq:MV}
\cdots\to\HH_q(A)\xrightarrow{(f_*,-g_*)}
\HH_q(B)\oplus\HH_q(C)\to\HH_q(B\hpo_A C)
\xrightarrow{\delta}\HH_{q-1}(A)\to\cdots.
\end{equation}

For maps of spans, let $u:A\to A'$, $v:B\to B'$, $w:C\to C'$ have homotopies $vf\sim f'u$ and $wg\sim g'u$. After cellular approximation, retain their chain homotopies $H,J$, with signs fixed by
\[
\partial H+H\partial=v_*f_*-f'_*u_*;
\qquad
\partial J+J\partial=w_*g_*-g'_*u_*.
\]
The induced chain map is
\[
(b,c,a)\longmapsto
\bigl(v_*b+Ha,\ w_*c-Ja,\ u_*a\bigr).
\]
Projection to the third coordinate explains the boundary-naturality formula. For finite CW spans, the formalized comparison with actual augmented cellular chains supplies $\delta_n:\HH_{n+1}(B\hpo_A C)\to\HH_n(A)$ and exactness at $\HH_{n+1}(B\hpo_A C)$, $\HH_n(A)$, and $\HH_n(B)\oplus\HH_n(C)$ for $0\leq n<5$. It supplies naturality for $\delta_2:H_3\to H_2$ between spans and for $\delta_1:H_2\to H_1$ under self-maps of a span, retaining the specified comparison homotopies. The presentations, cellular approximations, and marked homotopies are chosen jointly inside a single propositional truncation. These are the geometric interfaces used below; see Appendix~\ref{app:formal}.

\subsection{The component groups and covering degree}

Write $C\simeq S^1$ for the circle used in the Seifert construction. The boundary maps into $P_0\times C$ are $\partial_i\times\id_C$, where $\partial_1,\partial_2,\partial_3$ represent $a,b,(ab)^{-1}$. In the coordinates $a\times h,b\times h$, their degree-two columns are $(1,0),(0,1),(-1,-1)$. Hence the left map on $H_2$ is
\begin{equation}\label{eq:D}
D:\Z^3\to\Z^2,\qquad D(x,y,z)=(x-z,y-z),
\end{equation}
with section $(u,v)\mapsto(u,v,0)$ and kernel $\Z(1,1,1)$. The following lemma identifies the coefficients of these actual maps and of the overlap coverings.

\begin{lemma}[Torus multiplier]\label{lem:torus-multiplier}
For $\ell\in\Z$, the map $m_\ell(u,t)=(\mu_\ell(u),t)$ of $T^2=C\times C$ acts by multiplication by $\ell$ on $H_2$.
\end{lemma}
\begin{proof}
Use the cylinder presentation
\[
C\times\Sigma 2\simeq
C\hpo_{C\amalg C}C,
\]
whose two structure maps are folds and whose suspension meridians are the two pushout paths. Since $\Sigma2\simeq S^1$ and both vertices have zero $H_2$, its boundary $\delta:H_2(T^2)\to H_1(C\amalg C)$ is injective. The overlap map $\mu_\ell\amalg\mu_\ell$ acts by $\ell$ on $H_1$. Naturality for this one sequence gives $\delta H_2(m_\ell)(z)=\ell\delta(z)=\delta(\ell z)$, and injectivity gives the result.
\end{proof}

Coordinate projections select the corresponding cellular generators. Reflection in one circle reverses the orientation of the torus square cell and therefore has degree $-1$ on $H_2(T^2)$.

For positive $p,q,r$, use the overlap basis $(u_1,u_2,u_3,h_1,h_2,h_3)$ of $H_1(\coprod T^2)$, where $u_i$ is the boundary-circle loop and $h_i$ the fibre loop. In the target basis $(a,b,h,x,y,z)$, the boundary words $a,b,(ab)^{-1}$ and the filling rows $[1,p],[1,q],[1-r,r]$ give the signed Mayer--Vietoris map
\[
A_{\mathrm{raw}}=
\begin{pmatrix}
1&0&-1&0&0&0\\
0&1&-1&0&0&0\\
0&0&0&1&1&1\\
-1&0&0&-p&0&0\\
0&-1&0&0&-q&0\\
0&0&r-1&0&0&-r
\end{pmatrix}.
\]
Make the following integral changes of basis in its source and target, respectively:
\[
\begin{gathered}
(u_1,u_2,t_p,t_q,t_r,t_0)
=(u_1,u_2,-h_1,-h_2,-h_3,u_1+u_2+u_3+h_3),\\
(a-x,b-y,x,y,z,\bar h),\qquad \bar h=-h.
\end{gathered}
\]
Both are unimodular. The first two columns now form an identity block, and the remaining columns have zero coordinates in its two rows, giving $I_2\oplus A(p,q,r)$, with labelled rows and columns
\[
A(p,q,r)=
\begin{array}{c|rrrr}
 &t_p&t_q&t_r&t_0\\\hline
x&p&0&0&-1\\
y&0&q&0&-1\\
z&0&0&r&-1\\
\bar h&1&1&1&-1
\end{array}.
\]
Consequently,
\[
\det A=pq+pr+qr-pqr
=pqr\left(\frac1p+\frac1q+\frac1r-1\right).
\]
The second equality uses $p,q,r>0$. Surjectivity of $D$ and the adjacent Mayer--Vietoris terms now give $H_3\cong\ker D\cong\Z$ and $H_2\cong\ker A$. The degree-zero map is injective: the overlap and the filling circles both have reduced $H_0\cong\Z^2$, and the filling map preserves their component indices. Thus $H_1\cong\operatorname{coker}A$. The Smith diagonals and homology groups are
\begin{equation}\label{eq:component-homology}
\begin{array}{c|c|cccc}
X&\text{Smith diagonal of }A&H_1(X)&H_2(X)&H_3(X)&H_4(X)\\\hline
M(2,3,4)&(1,1,1,2)&\Z/2&0&\Z&0\\
M(2,3,5)&(1,1,1,1)&0&0&\Z&0\\
M(2,2,4)&(1,1,2,2)&(\Z/2)^2&0&\Z&0\\
M(2,2,3)&(1,1,1,4)&\Z/4&0&\Z&0
\end{array}.
\end{equation}
Each determinant is nonzero, so each kernel vanishes. These types have finite CW models of dimension $3$, and hence have zero homology in every higher degree.

\begin{lemma}[Boundary incidence]\label{lem:boundary-incidence}
Consider the initial lifted Seifert diagram, before the primitive cancellations. Let $Q_c$ be its sheet set and $\mathcal B_c$ the disjoint union of its three boundary-orbit sets. Orient the lifted graph edges as $e_a(s):s\to a(s)$ and $e_b(s):s\to b(s)$ for $s:Q_c$. Orient boundary loops by $a$, $b$, and $(ab)^{-1}$, with the last traversing $b^{-1}$ then $a^{-1}$, and orient each torus cell by this boundary direction followed by the fibre direction $h$. Write $A(s),B(s),C(s)$ for the three boundary orbits through $s$.

Choose a spanning tree in the lifted graph and let $J_c$ be its non-tree edges. The left $H_2$ map is the matrix $\widehat D_c:\Z[\mathcal B_c]\to\Z[J_c]$ obtained by restricting the following rows to $J_c$:
\[
\begin{aligned}
(R_cx)_{e_a(s)}&=x_{A(s)}-x_{C(s)},\\
(R_cx)_{e_b(s)}&=x_{B(s)}-x_{C(b(s))}.
\end{aligned}
\]
Thus each row of $\widehat D_c$ has a $+1$ at the indicated $A$- or $B$-orbit and a $-1$ at the indicated $C$-orbit, with all other entries zero. Its sizes and kernels are
\[
\begin{array}{c|ccc}
c&|Q_c|&|J_c|&|\mathcal B_c|\\\hline
O&3&4&5\\
I&10&11&12
\end{array},
\qquad
\ker\widehat D_c=\Z\mathbf1.
\]
\end{lemma}
\begin{proof}
Every positive $a$-edge occurs once in its $A$-boundary and once, negatively, in the $C$-boundary through its initial sheet. The analogous negative occurrence of $e_b(s)$ belongs to the $C$-boundary through $b(s)$. These occurrences give the displayed rows. Tree contraction reads off the non-tree coefficients, and the torus multiplier calculation identifies them with the actual $H_2$ coefficients. There are $2N_c$ edges and $N_c-1$ tree edges, so $|J_c|=N_c+1$. The orbit lists in \eqref{eq:orbit-lengths} give $|\mathcal B_O|=5$ and $|\mathcal B_I|=12$.

For any $x$, the chain $R_cx$ is a sum of closed boundary words and hence a graph cycle. If $\widehat D_cx=0$, this cycle is supported on the tree; successive removal of leaves gives $R_cx=0$. Put $\alpha(s)=x_{A(s)}$, $\beta(s)=x_{B(s)}$, and $\gamma(s)=x_{C(s)}$. The two zero-row equations give $\alpha(s)=\gamma(s)$ and $\beta(s)=\gamma(b(s))$. Since $\beta$ is constant on $b$-orbits, the latter equation also gives $\beta(s)=\gamma(s)$. Constancy of $\alpha$ on $a$-orbits and of $\beta$ on $b$-orbits makes $\gamma$ constant along every graph edge. The lifted graph is connected, so all boundary coefficients are equal. Conversely the displayed rows kill every constant vector. Evaluation at the $A$-orbit of a fixed root sheet is inverse to $t\mapsto t\mathbf1$, proving the asserted integral kernel identification.
\end{proof}

\begin{theorem}[Covering degree]\label{thm:cover-degree}
There are jointly chosen coordinates $H_3(A_c)\cong\Z$ and $H_3(C_c)\cong\Z$ in which the specified coverings $p_O,p_I$ act by $+3,+10$, respectively.
\end{theorem}
\begin{proof}
In the initial covering pushout, let $O_s,O_t$ be the source and target overlaps, and let $V_s=\Z[\mathcal B_c]$, $W_s=\Z[J_c]$. The left $H_2$ maps have matrices $\widehat D=\widehat D_c:V_s\to W_s$ and $D:\Z^3\to\Z^2$. Lemma~\ref{lem:boundary-incidence} identifies the actual source map, with its $4\times5$ or $11\times12$ incidence matrix, and gives
\[
\ker\widehat D=\Z\mathbf1,\qquad \ker D=\Z(1,1,1).
\]
The overlap projection $\pi:O_s\to O_t$ is represented by $U_2:V_s\to\Z^3$, with coefficients the boundary-orbit lengths, by Lemma~\ref{lem:torus-multiplier}. For each boundary kind these lengths sum to the number $N_c$ of sheets; thus $U_2(\mathbf1)=N_c(1,1,1)$.

The vertices have zero $H_3$, and the right vertices have zero $H_2$, so the boundary maps identify source and target $H_3$ with the two kernels above. Naturality gives
\[
\delta_t H_3(p_c)=H_2(\pi)\delta_s.
\]
Choosing the $H_3$ generators by $\delta_s^{-1}(\mathbf1)$ and $\delta_t^{-1}(1,1,1)$ proves the asserted positive degrees. The specified equivalences from the reduced covering diagrams transport these coordinates and their equation to $A_c,C_c$.
\end{proof}

\subsection{The outer pushouts and local contraction}

For $E_c=R\hpo_{A_c}C_c$, the groups $H_2(R),H_2(A_c),H_2(C_c)$ vanish. Theorem~\ref{thm:cover-degree} gives the exact segment
\[
0\longrightarrow H_4(E_c)\longrightarrow\Z
 \xrightarrow{\times N_c}\Z\longrightarrow H_3(E_c)\longrightarrow0,
\]
where the sign in \eqref{eq:MV} is absorbed in the target coordinate. Consequently $H_4(E_c)=0$ and $H_3(E_c)\cong\Z/N_c$. The next segment identifies
\[
H_2(E_c)\cong\ker\bigl(H_1(A_c)\xrightarrow{(q_{c*},-p_{c*})}
H_1(R)\oplus H_1(C_c)\bigr).
\]
By Proposition~\ref{prop:seifert-simple}, $H_1(E_c)=0$, so this map is surjective. For $c=O$ it is a surjection $(\Z/2)^2\to(\Z/2)^2$, hence an isomorphism. For $c=I$ it is a surjection $\Z/4\to\Z/2$, with kernel $\Z/2$.

\begin{proposition}[Homology profiles]\label{prop:profiles}
The types $E_O,E_I$ are simply connected and merely admit finite CW models of dimension at most $4$. Their nonzero reduced homology groups are
\[
\HH_3(E_O)\cong\Z/3;\qquad
\HH_2(E_I)\cong\Z/2,\quad \HH_3(E_I)\cong\Z/10.
\]
\end{proposition}
\begin{proof}
The homology calculation is above. The double mapping cylinder for each outer span has dimension at most $4$, since $R$ has dimension $2$ and $A_c,C_c$ have dimension $3$.
\end{proof}

For $E_O$, Hurewicz and $H_2(E_O)=0$ give $2$-connectedness. Theorem~\ref{thm:moore-recognition} yields $\mere{E_O\simeq\Moore33}$. Eliminating this truncation into the proposition $\isContr(L_3E_O)$ and using \eqref{eq:needed-moore-null} contracts $L_3E_O$.

The type $E_I$ has $H_2(E_I)\cong\Z/2$; we cone off $\Moore22$, leaving a $2$-connected cofibre to which Moore recognition applies. Choosing a point under truncation, Hurewicz in degree two supplies, merely, a representative $s:S^2\to_*E_I$ of the generator. Since $H_2(\delta_{2,2})=\times2$, naturality and injectivity of Hurewicz give, merely, a pointed nullhomotopy of $s\delta_{2,2}$. Cofibre recursion gives $a:\Moore22\to E_I$ inducing an isomorphism on $H_2$, calibrated on the bottom inclusion as in Step~1 above. Its cofibre $K_a$ is simply connected and has a finite CW model of dimension at most $4$, since $\Moore22$ has dimension $3$. The cofibre sequence gives
\[
H_2(K_a)=0,\qquad H_3(K_a)\cong\Z/10,\qquad H_4(K_a)=0.
\]
Hurewicz makes $K_a$ $2$-connected, and Theorem~\ref{thm:moore-recognition} gives the data in the form
\begin{equation}\label{eq:EI-profile}
\mere{\sum_{a:\Moore22\to E_I}
             \bigl(\cofib(a)\simeq\Moore{10}3\bigr)}.
\end{equation}
For each such $a$, both its domain and its cofibre are $L_{10}$-connected by \eqref{eq:needed-moore-null}. Lemma~\ref{lem:two-stage-null}, followed by elimination into the proposition $\isContr(L_{10}E_I)$, gives the second contraction.

\begin{corollary}\label{cor:local-totals}
The total spaces satisfy $L_3E_O\simeq1$ and $L_{10}E_I\simeq1$.
\end{corollary}
These contractions are used in \S\ref{sec:construction} to identify the localized fibres with loop types.

\section{Proof of the delooping assertion of Theorem~\ref{thm:main}}\label{sec:construction}

The inputs are the fibre sequences $F\to E_c\xrightarrow{t_c}P_c$ of \S\ref{sec:truncation}, with simply connected bases and specified common fibre $F\simeq S^3$, together with the contractions $L_3E_O\simeq1$ and $L_{10}E_I\simeq1$ from \S\ref{sec:homology}. Fibre preservation and arithmetic fracture give a pointed connected type $B$ and a pointed equivalence $\Omega B\eqp S^3$.

\subsection{Localization conventions}\label{sec:local-tools}
For $m\geq1$, write $L_m=L_{d_m}$ for localization inverting $m$, where $d_m:S^1\to S^1$ has degree $m$, and write $\eta^m_X:X\to L_mX$ for its unit. This is the $\mathsf{deg}(S)$-localization for the constant sequence $S(k)=m$. A type is $m$-local exactly when $p\mapsto p^m$ is an equivalence on each loop type \cite[Lemma~4.2]{cors}. A map $g:A\to D$ is an \emph{$L_m$-equivalence} if $L_mg$ is an equivalence; equivalently, precomposition $g^*:\Map(D,Z)\to\Map(A,Z)$ is an equivalence for every $m$-local $Z$ \cite[Lemma~2.9]{cors}. A type $A$ is \emph{$L_m$-connected} if $L_mA\simeq1$.

For $m$-local $Z$, restriction along $\eta^m_X$ gives $\Map(L_mX,Z)\simeq\Map(X,Z)$; taking fibres of evaluation gives the pointed version. Local types are closed under identity types and fibres of maps between local types \cite[\S2.1]{cors}, and $L_m$ preserves connectedness in every degree \cite[Corollary~4.14]{cors}. The local Moore contractions and two-stage nullity argument are Corollary~\ref{cor:moore-null} and Lemma~\ref{lem:two-stage-null}. Their application in Corollary~\ref{cor:local-totals} gives
\begin{equation}\label{eq:local-totals}
L_3E_O\simeq1,\qquad L_{10}E_I\simeq1.
\end{equation}

\subsection{Fibre preservation}\label{sec:fibres}
The following is the fibre-preservation statement in \cite[\S4, observation following Proposition~4.10]{scoccola}, specialized to the constant degree family. The earlier \cite[Lemma~4.16]{cors} assumes that both the base and total space are simply connected; Scoccola's statement requires only the base.
\begin{theorem}[Fibre preservation over a simply connected base]\label{thm:local-fibre}
Let $f:E\to Y$, where $Y$ is simply connected. For every $y:Y$ the canonical comparison
\begin{equation}\label{eq:local-fibre}
\lambda_{f,y}:L_m\fib(f,y)\xrightarrow{\simeq}\fib(L_mf,\eta^m_Yy)
\end{equation}
is an equivalence. There is no connectedness assumption on $E$ or on the fibre.
\end{theorem}
\begin{proof}
Put $L'_m=L_{\Sigma d_m}$, with unit $\eta'^m_X:X\to L'_mX$. This is the separated reflection: a type is $L'_m$-local precisely when its identity types are $m$-local \cite[Lemma~2.15]{cors}. Proposition~4.7 of \cite{scoccola} extends the family $y\mapsto L_m\fib(f,y)$ to an $m$-local-valued family over $L'_mY$, with total space $E'$ and an $L_m$-equivalence $E\to E'$. Since $Y$ is simply connected, Proposition~4.3 of \cite{scoccola} identifies $L'_mY$ with $L_mY$ compatibly with units. This base remains simply connected by \cite[Corollary~4.14]{cors}. Lemma~4.9 of \cite{scoccola} makes $E'$ local, so $E\to E'$ identifies it with $L_mE$ over $L_mY$. Its fibres give the asserted equivalence.
\end{proof}
For later compatibility we specify the comparison. If $\nu_f(e):L_mf(\eta^m_Ee)=\eta^m_Y(fe)$ is unit naturality, $\lambda_{f,y}$ is the unique extension into the local target in \eqref{eq:local-fibre} of
\[
(e,p)\longmapsto\bigl(\eta^m_Ee,\nu_f(e)\cdot\operatorname{ap}_{\eta^m_Y}(p)\bigr).
\]
Thus this formula is its computation on units, with the specified naturality path.

\begin{corollary}[Local delooping machine]\label{cor:local-machine}
Suppose $f:E\to Y$ has points $e_0,y_0$ and a specified path $s:f(e_0)=y_0$. If $Y$ is simply connected and $L_mE$ is contractible, then
\[
\Omega(L_mY,\eta^m_Yy_0)\simeq L_m\fib(f,y_0).
\]
\end{corollary}
\begin{proof}
Use \eqref{eq:local-fibre}, centre the contraction at $\eta^m_Ee_0$, and identify its fibre with loops using $\nu_f(e_0)\cdot\operatorname{ap}_{\eta^m_Y}(s)$. This retains the specified pointing path.
\end{proof}

\subsection{Arithmetic fracture}\label{sec:fracture}
For positive integers $m,n$, the type $L_{mn}X$ is both $m$-local and $n$-local. The localization mapping property therefore supplies maps and specified homotopies
\[
\begin{aligned}
j_m &: L_mX\longrightarrow L_{mn}X,
& h_m &: j_m\eta^m_X\sim\eta^{mn}_X,\\
j_n &: L_nX\longrightarrow L_{mn}X,
& h_n &: j_n\eta^n_X\sim\eta^{mn}_X.
\end{aligned}
\]
\begin{theorem}[Arithmetic fracture {\cite[Theorem~5.4]{scoccola}}]\label{thm:fracture}
Let $X$ be simply connected and $m,n\geq1$. If $am+bn=1$ for integers $a,b$, then the map
\begin{equation}\label{eq:fracture-map}
\begin{aligned}
\gamma_X &: X\longrightarrow L_mX\hpb_{L_{mn}X}L_nX,\\
\gamma_X(x)&=\bigl(\eta^m_Xx,\eta^n_Xx,h_m(x)\cdot h_n(x)^{-1}\bigr)
\end{aligned}
\end{equation}
is an equivalence.
\end{theorem}
To specialize the cited theorem, take the constant sequences $R_k=n$, $S_k=m$ and their pointwise product $T_k=mn$. Its cross-coprimality hypothesis says that every $R_k$ is coprime to every $S_l$, here exactly $\gcd(m,n)=1$. Lemma~5.1 of \cite{scoccola} identifies the composites through $L_mX$ and $L_nX$ as $L_{mn}$-localizations, giving unit-compatible comparisons with the iterated localizations in that theorem. Thus its pullback map is \eqref{eq:fracture-map}. The formalization proves the $2$-connected case, which covers the instance $X=F\simeq S^3$ used here.

For $m=3$, $n=10$, take $a=7$ and $b=-2$, so $3\cdot7+10\cdot(-2)=1$. The resulting fracture square is
\begin{equation}\label{eq:sphere-fracture}
\begin{tikzcd}[column sep=large,row sep=large]
S^3 \arrow[r,"\eta^3_{S^3}"] \arrow[d,"\eta^{10}_{S^3}"'] \pbmark & L_3S^3 \arrow[d,"j_3"] \\
L_{10}S^3 \arrow[r,"j_{10}"'] & L_{30}S^3.
\end{tikzcd}
\end{equation}

\subsection{The common comparison}\label{sec:common}
Keep $R=\RP^2$, $\baseN=\trunc2R$, and the specified common fibre $F\simeq S^3$ of \S\ref{sec:truncation}. Form
\[
\begin{tikzcd}[column sep=large,row sep=large]
R \arrow[r] \arrow[d] \pomark & E_I \arrow[d] \\
E_O \arrow[r] & E_*
\end{tikzcd}
\qquad
\begin{tikzcd}[column sep=large,row sep=large]
\baseN \arrow[r] \arrow[d] \pomark & P_I \arrow[d] \\
P_O \arrow[r] & P_*.
\end{tikzcd}
\]
Cartesian descent applied to the two squares $R\to E_c$ over $\baseN\to P_c$ gives $F\to E_*\xrightarrow{t_*}P_*$. The induced maps from the $O$ and $I$ fibre sequences are the identity in the specified $F$ coordinate. Moreover, $P_*$ is simply connected: $P_O,P_I$ are simply connected, $\baseN$ is connected, and the pushout van Kampen theorem gives a contractible $1$-truncation.

The cofibre presentation $R=\RP^2\simeq\Moore21$ and Corollary~\ref{cor:moore-null} give $L_{30}R\simeq1$. Every $30$-local type $Z$ is $3$-local and $10$-local, since the corresponding commuting power maps compose to the thirtieth-power equivalence. Together with \eqref{eq:local-totals}, these facts identify all three constant-map inclusions into $\Map(E_O,Z)$, $\Map(E_I,Z)$ and $\Map(R,Z)$ as equivalences. The pushout mapping property gives
\begin{align*}
\Map(E_*,Z)
&\simeq\Map(E_O,Z)\hpb_{\Map(R,Z)}\Map(E_I,Z)\\
&\simeq Z\hpb_ZZ\simeq Z.
\end{align*}
These equivalences identify the constant-map inclusion, so $L_{30}E_*\simeq1$.

Put $Q_O=L_3P_O$, $Q_I=L_{10}P_I$ and $T=L_{30}P_*$, pointed by the units. Let $u_O:Q_O\to T$ and $u_I:Q_I\to T$ be induced by the pushout inclusions, with their pointing paths. Corollary~\ref{cor:local-machine} gives
\[
\phi_O:\Omega Q_O\simeq L_3F,\qquad
\phi_I:\Omega Q_I\simeq L_{10}F,\qquad
\phi_*:\Omega T\simeq L_{30}F.
\]
They fit into specified commuting squares
\begin{equation}\label{eq:loop-comparisons}
\begin{tikzcd}[column sep=large,row sep=large]
\Omega Q_O \arrow[r,"\Omega u_O"] \arrow[d,"\phi_O"'] & \Omega T \arrow[d,"\phi_*"] \\
L_3F \arrow[r] & L_{30}F
\end{tikzcd}
\qquad
\begin{tikzcd}[column sep=large,row sep=large]
\Omega Q_I \arrow[r,"\Omega u_I"] \arrow[d,"\phi_I"'] & \Omega T \arrow[d,"\phi_*"] \\
L_{10}F \arrow[r] & L_{30}F.
\end{tikzcd}
\end{equation}
To obtain each homotopy, precompose with $\phi_c^{-1}$ and restrict along $F\to L_{N_c}F$. The common-fibre identifications and the unit rule in Theorem~\ref{thm:local-fibre} give the equality there. Contractions centred at the chosen points are compatible because the type of such contractions is contractible. Restriction into $L_{30}F$ is an equivalence of function types, so it extends this equality with its prescribed restriction. The pointed version retains the basepoint paths. In particular, the $I$-side point is compared with the selected star point by the pushout gluing path.

\subsection{Assembly}\label{sec:assembly}
\begin{proof}[Proof of the delooping assertion of Theorem~\ref{thm:main}]
The type in \eqref{eq:B} is explicitly
\[
B=\sum_{x:Q_O}\sum_{y:Q_I}(u_Ox=u_Iy).
\]
Write $q_O^0,q_I^0,t_0$ for the chosen points and $\alpha_c:u_c(q_c^0)=t_0$ for the pointing paths. Its basepoint is
\[
b_0=(q_O^0,q_I^0,\alpha_O\cdot\alpha_I^{-1}).
\]
The types $Q_O,Q_I,T$ are simply connected. To compare two points of $B$ merely, use connectedness of $Q_O,Q_I$ to compare their first two coordinates. After transport, the last coordinates are paths in $T$ with fixed endpoints; their identity type is merely inhabited by simple connectedness of $T$. Elimination into the proposition expressing mere equality gives $c_B:\Conn(B)$.

Looping the pointed homotopy pullback and using \eqref{eq:loop-comparisons} gives the chain recorded in \eqref{eq:loop-chain}:
\begin{align*}
\Omega(B,b_0)
&\simeq\Omega Q_O\hpb_{\Omega T}\Omega Q_I\\
&\simeq L_3F\hpb_{L_{30}F}L_{10}F
 \simeq F\simeq S^3.
\end{align*}
The fracture equivalence applies to the simply connected type $F\simeq S^3$. Let $e:\Omega(B,b_0)\simeq S^3$ be this composite and put $r=e^{-1}(*)$. Left translation $\ell_r(p)=r\cdot p$ is an equivalence, with inverse $\ell_{r^{-1}}$. The composite $e\ell_r$ has the pointing
\[
e(r\cdot\mathsf{refl})=e(r)=*,
\]
by the right unit law and the inverse law of $e$. It defines $\varepsilon:\Omega(B,b_0)\eqp S^3$, and hence the required untruncated term
\[
\mathfrak d=((B,b_0),c_B,\varepsilon):
\sum_{B:\Pointed}\bigl(\Conn(B)\times(\Omega B\eqp S^3)\bigr).
\]
Homology coordinates and Moore representatives enter only into proofs of contractibility, which are propositions. Their propositional truncations can therefore be eliminated there; the pullback $B$ and the displayed term are then constructed from the fixed geometric data.
\end{proof}

\section{The maximal torus and comparison problems}\label{sec:outlook}

The common corner of the construction yields more than a delooping: we will identify $R=\RP^2$ as the fibre of a map $\baseN\to B$. Its orientation cover then gives a map $K(\Z,2)\to B$ with fibre $S^2$. This maximal torus identifies the classical homotopy type of $B$; internal comparisons with standard structures are discussed below.

\subsection{The maximal torus and classical recognition}

Put $u=\eta_2:R\to\baseN$, $n_0=u(r_0)$ and $E=L_3E_O\hpb_{L_{30}E_*}L_{10}E_I$. The three localized total spaces are contractible, so $E$ is contractible. The common corner induces a square
\begin{equation}\label{diag:HP}
\begin{tikzcd}[column sep=large,row sep=large]
R \arrow[r,"\alpha"] \arrow[d,"u"'] \pbmark & E \arrow[d,"p"] \\
\baseN \arrow[r,"j"'] & B.
\end{tikzcd}
\end{equation}
Our goal is its cartesianness. Contracting $E$ will then identify $R$ with $\fib(j,b_0)$ over $\baseN$, allowing us to pull back the orientation cover. The following lemma isolates the arithmetic step, including the specified overlap path. We use localization HIT recursors with judgmental computation on units; computation paths for homotopies between recursors need only be propositional.

\begin{lemma}[Recognition of a specified fracture cone]\label{lem:unit-fracture-cone}
Let $F$ be $2$-connected and $m,n$ coprime positive integers. Let $X_i$ be $i$-local for $i=m,n,mn$, with maps $a:X_m\to X_{mn}\leftarrow X_n:b$. Suppose
$c=(c_m,c_n,\gamma):F\to X_m\hpb_{X_{mn}}X_n$ is a specified cone. If the recursor extensions $e_m:L_mF\to X_m$, $e_n:L_nF\to X_n$ and $e_{mn}:L_{mn}F\to X_{mn}$ of $c_m,c_n,a c_m$, respectively, are equivalences, then $c$ is an equivalence.
\end{lemma}
\begin{proof}
Choose the overlap maps $j_i:L_iF\to L_{mn}F$ by recursion, so $j_i\eta^i\equiv\eta^{mn}$, and choose their unit homotopies to be reflexivity. Thus the fracture map is $\gamma_F(z)=(\eta^mz,\eta^nz,\mathsf{refl})$. Since $X_{mn}$ is both $m$-local and $n$-local, restriction along the units is an equivalence of mapping spaces and of their identity types. It therefore supplies homotopies with specified propositional unit restrictions
\[
\begin{aligned}
A:a e_m\sim e_{mn}j_m,&\qquad A(\eta^mz)=\mathsf{refl},\\
D:b e_n\sim e_{mn}j_n,&\qquad D(\eta^nz)=\gamma(z)^{-1}.
\end{aligned}
\]
The objectwise equivalences $e_i$ induce the pullback equivalence
\[
\Phi(x,y,s)=\bigl(e_mx,e_ny,
 A(x)\cdot\operatorname{ap}_{e_{mn}}(s)\cdot D(y)^{-1}\bigr).
\]
On $\gamma_F(z)$ its first two coordinates are $c_mz,c_nz$; the last is, by the specified restrictions and path laws,
$\mathsf{refl}\cdot\mathsf{refl}\cdot(\gamma(z)^{-1})^{-1}=\gamma(z)$.
Thus $\Phi\gamma_F\sim c$, including the actual overlap coordinate. Theorem~\ref{thm:fracture} makes $\gamma_F$, and hence $c$, an equivalence.
\end{proof}

\begin{theorem}\label{thm:torus}
The square \eqref{diag:HP} is cartesian, and there is a pointed map $t:K(\Z,2)\to_*B$ with $\fib(t)\simeq S^2$. This is formalized in safe Cubical Agda (Appendix~\ref{app:formal}).
\end{theorem}
\begin{proof}
\emph{The specified square.}
Write $i_c:R\to E_c$, $j_c:\baseN\to P_c$ for the common-corner inclusions and $k_c^E:E_c\to E_*$, $k_c^P:P_c\to P_*$ for the star inclusions ($c=O,I$). The structure maps $t_c,t_*$ satisfy $t_ci_c=j_cu$ and $t_*k_c^E=k_c^Pt_c$ judgmentally. The actual pushout seams satisfy
\[
\begin{gathered}
\sigma_R(r):k_O^Ei_O(r)=k_I^Ei_I(r),\qquad
\sigma_N(n):k_O^Pj_O(n)=k_I^Pj_I(n),\\
\operatorname{ap}_{t_*}(\sigma_R(r))=\sigma_N(u r).
\end{gathered}
\]
Put $V_c=L_{N_c}E_c$, $V_*=L_{30}E_*$ and denote the localized maps by $v_c:V_c\to V_*$, $t_c^L:V_c\to Q_c$ and $t_*^L:V_*\to T$. All these maps, and $u_c:Q_c\to T$, compute on units by their original maps. Restriction into the local type $T$ extends reflexivity on units to side homotopies with specified computations
\[
H_c(x):t_*^L(v_cx)=u_c(t_c^Lx),\qquad
\beta_c(x):H_c(\eta^{N_c}x)=\mathsf{refl}.
\]
The $\beta_c$ are propositional paths, not judgmental rules. Define
\[
\begin{aligned}
p(x,y,s)&=\bigl(t_O^Lx,t_I^Ly,
 H_O(x)^{-1}\cdot\operatorname{ap}_{t_*^L}(s)\cdot H_I(y)\bigr),\\
\alpha(r)&=\bigl(\eta^3 i_O(r),\eta^{10}i_I(r),
                    \operatorname{ap}_{\eta^{30}}(\sigma_R(r))\bigr),\\
j(n)&=\bigl(\eta^3 j_O(n),\eta^{10}j_I(n),
                    \operatorname{ap}_{\eta^{30}}(\sigma_N(n))\bigr).
\end{aligned}
\]
The seam identity and $\beta_O,\beta_I$ give $H_r:p(\alpha(r))=j(u r)$, with constant first two coordinates. The specified pointings identify the seam of $j(n_0)$ with that of $b_0$, giving $j_0:j(n_0)=b_0$.

\emph{The fibre comparison.}
The square induces
\[
\kappa_n:\fib(u,n)\longrightarrow\fib(p,j(n)),\qquad
\kappa_n(r,q)=\bigl(\alpha(r),H_r\cdot\operatorname{ap}_j(q)\bigr).
\]
At $n_0$, put $b_c^L=\eta^{N_c}j_c(n_0)$ and $\sigma=\operatorname{ap}_{\eta^{30}}(\sigma_N(n_0))$. The standard fibre formula for a map of pullbacks gives
\[
\begin{gathered}
\Psi:\fib(p,j(n_0))\simeq X_3\hpb_{X_{30}}X_{10},\\
X_3=\fib(t_O^L,b_O^L),\quad X_{10}=\fib(t_I^L,b_I^L),\quad
X_{30}=\fib(t_*^L,u_Ob_O^L),\\
a(x,q)=(v_Ox,H_O(x)\cdot\operatorname{ap}_{u_O}(q)),\qquad
b(y,q)=(v_Iy,H_I(y)\cdot\operatorname{ap}_{u_I}(q)\cdot\sigma^{-1}).
\end{gathered}
\]
Indeed, the overlap square in a path $p(x,y,s)=j(n_0)$ becomes a path $a(x,q_O)=b(y,q_I)$ over $s$ after whiskering by $H_O,H_I,\sigma^{-1}$; reversing the whiskerings gives the inverse. This description preserves the fibre-coordinate projections.

Set $F=\fib(u,n_0)$ and write the actual composite as $\Psi\kappa_{n_0}=(c_3,c_{10},\gamma)$. Its coordinate maps on $(r,q):F$ are, using the constant coordinates of $H_r$ and the unit laws,
\[
\begin{aligned}
c_{N_c}(r,q)&=\bigl(\eta^{N_c}i_c(r),
       \operatorname{ap}_{\eta^{N_c}j_c}(q)\bigr),\\
a c_3(r,q)&=\bigl(\eta^{30}k_O^Ei_O(r),
       \operatorname{ap}_{\eta^{30}k_O^Pj_O}(q)\bigr).
\end{aligned}
\]
These are pointwise paths, not judgmental rules; the second uses $\beta_O$. These are exactly the common-fibre equivalences supplied by cartesian descent followed by the fibre-unit maps of Theorem~\ref{thm:local-fibre}. Each $X_i$ is $i$-local, and the original bases $P_O,P_I,P_*$ are simply connected. Fibre preservation therefore makes all three recursor extensions equivalences: their unit restrictions are the displayed maps, so the localization mapping property identifies them with the canonical fibre equivalences. Since $F\simeq S^3$ is $2$-connected, Lemma~\ref{lem:unit-fracture-cone} applies to this actual cone, proving $\kappa_{n_0}$ an equivalence without changing its seam $\gamma$.

The type $\baseN$ is connected and $\isEquiv(\kappa_n)$ is a proposition, so transport from $n_0$, eliminating mere paths, proves it for every $n$. Taking dependent sums gives the equivalence
$R\simeq\sum_{n:\baseN}\fib(p,j(n))\simeq E\hpb_B\baseN$;
on $r$ its value is $(\alpha(r),u r,H_r)$, up to the unit law. Thus the specified square is cartesian.

\emph{The pointed torus.}
Contract $E$ at $e_0=\alpha(r_0)$, writing $\ell_e:e=e_0$, and put $p_0=H_{r_0}\cdot j_0$. Rotating the pullback gives an equivalence \emph{over $\baseN$}
\[
\chi:R\simeq\fib(j,b_0),\qquad
\chi(r)=\bigl(u r,H_r^{-1}\cdot\operatorname{ap}_p(\ell_{\alpha(r)})\cdot p_0\bigr).
\]
For the orientation cover $q:S^2\to R$, put $T_2=\tau_2S^2$ and $a_2=\tau_2q:T_2\to\baseN$. Its set fibres make its level-$2$ truncation square cartesian by Theorem~\ref{thm:trunc-square}. Pulling back $\chi$ along $a_2$ therefore gives
\[
S^2\simeq\sum_{z:T_2}\fib(u,a_2z)
 \simeq\sum_{z:T_2}(j(a_2z)=b_0)=\fib(j a_2,b_0).
\]
Lemma~\ref{lem:hopf-reflection} supplies a pointed equivalence $e:T_2\eqp K(\Z,2)$. Define $t=j a_2e^{-1}$. Its pointing is the composite induced by $e^{-1}(0)=\eta_2(*)$, $a_2(\eta_2(*))=n_0$ and $j_0$. Change of domain along $e^{-1}$ now gives $\fib(t,b_0)\simeq S^2$.
\end{proof}

In the classical homotopy interpretation, $K(\Z,2)\simeq\CP^\infty=\B S^1$, and the map $t$ has finite homotopy fibre $S^2$. Thus it is a maximal torus for the delooping $B$.

\begin{corollary}[Classical recognition]
The classical homotopy type of $B$ is equivalent to $\HP^\infty$.
\end{corollary}
\begin{proof}
By Dwyer--Miller--Wilkerson \cite[Thm.~1.1 and Rem.~1.2]{dmw}, all deloopings of $S^3$ belong to the genus classified by Rector. The equivalence $\Omega B\eqp S^3$ places $B$ in this class. Rector and McGibbon's maximal-torus criterion then identifies $B$ with $\B S^3\simeq\HP^\infty$ \cite{rector,mcgibbon}; see also \cite[Thm.~1.1(4)]{yaumaps}. Its finite-fibre hypothesis is supplied by Theorem~\ref{thm:torus}.
\end{proof}
This corollary uses classical recognition results. An internal equivalence between $B$ and an independently defined $\HP^\infty$ remains open.

\subsection{Comparison with the Cayley--Dickson multiplication}

The Cayley--Dickson multiplication $\mu_{\mathrm{CD}}$ on $S^3$ \cite{cayley} is an $H$-space structure, and the delooping gives a second multiplication $\mu_B$ on $S^3$, $\mu_B(x,y)=\varepsilon(\varepsilon^{-1}(x)\cdot\varepsilon^{-1}(y))$, with $\varepsilon$ the pointed equivalence of Theorem~\ref{thm:main}.

\begin{problem}\label{prob:CD}
Construct a pointed $H$-space equivalence
\[
h:(S^3,\mu_B)\simeq_H(S^3,\mu_{\mathrm{CD}}).
\]
Thus $h:S^3\eqp S^3$ should admit a homotopy
\[
h(\mu_B(x,y))=\mu_{\mathrm{CD}}(h(x),h(y))
\qquad(x,y:S^3)
\]
compatible with the left and right unit homotopies.
\end{problem}

The composite $h\circ\varepsilon$ would then identify $B$ as a delooping of the Cayley--Dickson $H$-space. This formulation allows the sphere identification to vary. The quaternionic Hopf sequence $S^3\to S^7\to S^4$ constructed in \cite{cayley} is distinct from the complex Hopf sequence $S^1\to S^3\to S^2$ used in \S\ref{sec:truncation}. Relating these structures to the maximal torus of Theorem~\ref{thm:torus} requires further compatibility data.

\subsection{Further perspectives}

The construction illustrates how localization and arithmetic fracture, developed internally in homotopy type theory \cite{cors,scoccola}, can be used to construct a classical homotopy-theoretic object without first reproducing its customary geometric model. Here the higher coherence problem is approached through local constructions and arithmetic reconstruction, rather than by specifying an associator and then successively filling higher coherence diagrams.

This suggests a broader direction for synthetic homotopy theory: extending its local and arithmetic methods, and using them to construct and study further classical objects internally. Natural longer-term goals include synthetic treatments of Bott periodicity and the further development of synthetic stable homotopy theory. I hope that this example will encourage further work along these lines.

\appendix
\section{Formalization}\label{app:formal}

\subsection{Statements and scope}
The development uses Cubical Agda \cite{cubicalagda} and the Cubical library \cite{cubical}. The entry \path{agda/src/DeloopS3/Assembly/Complete.agda} exports
\[
\texttt{Deloop-S3}:
\sum_{B:\mathcal U_{0,*}}
  \bigl(\Conn(B)\times(\Omega B\eqp S^3)\bigr).
\]
Here $\mathcal U_{0,*}$ is \texttt{Pointed} at level zero;
$\Conn(B)$ is \texttt{isConnected 2 (typ B)}, and the pointed sphere
is written $(\texttt{S}_{+}\;3,\mathsf{north})$ in the source.
The sum is untruncated. The term has no module parameters and uses no
postulates. The second entry, \path{agda/src/DeloopS3/HPInf/Complete.agda},
exports the pointed map \texttt{t}, its fibre equivalence \texttt{t-fiber},
and the combined declaration
\texttt{S\textsuperscript{3}-delooping-with-maximal-torus}.
The latter uses the same concrete $B$: its map is $K(\Z,2)\to_*B$
and its fibre over the basepoint is equivalent to $S^2$. It does not assert
an equivalence with an independently defined $\mathbb H P^\infty$.
The pointed loop equivalence is obtained by translating the source loop
space before applying the unpointed equivalence
(\texttt{point-loop-equivalence} in \path{Assembly/DefB.agda}); it does
not assert that the unmodified comparison preserves the basepoint or that
it respects the Cayley--Dickson multiplication.
Both entries use \texttt{--safe --cubical --guardedness}.
Some auxiliary modules use \texttt{--lossy-unification}, an elaboration
option with no additional axioms.

The formalized fracture theorem assumes $2$-connectedness, which suffices
for $F\simeq S^3$; the simply connected theorem cited in the paper is the
more general result of \cite[Thm.~5.4]{scoccola}. Moore recognition is
formalized for arbitrary positive $m$ and $2$-connected finite CW types
of dimension at most $4$ with the stated homology profile. The bottom-Moore
cofibre argument and the resulting contractions are formalized for the
$O$ and $I$ profiles used here. Classical recognition uses
\cite[Rem.~1.2]{dmw}, without a separate genus-membership assumption;
it is cited, not formalized. Independent $\mathbb H P^\infty$ and Cayley--Dickson
multiplication comparisons remain outside the exports.

The geometric Mayer--Vietoris interface supplies exactness for $0\leq n<5$,
with boundary $H_{n+1}\to H_n$. Its naturality comparisons establish
$H_3\to H_2$ naturality for maps between spans and $H_2\to H_1$ naturality
for self-span maps, retaining their comparison homotopies. These are the
degrees used in \S\ref{sec:homology}. The all-degree chain equivalence
and these degree-specific geometric comparisons have separate exports.

\medskip
\noindent\textbf{Artifact.} The formalization and accompanying paper are hosted at \url{https://github.com/wxkj99/s3-delooping}. The formalization source snapshot used for this paper is commit \href{https://github.com/wxkj99/s3-delooping/tree/2eda93321fa89daa4875e6921db3b5fae19f75f2}{\texttt{2eda933}}. The repository contains the formalization, checking tools, pinned Cubical dependency, and the paper's \LaTeX{} sources and PDF. Local audit documents and verification records are maintained separately.

\subsection{Indexing conventions}
The paper uses classical truncation degrees. The library uses the following
indices.
\begin{center}
\small
\begin{tabular}{@{}ll@{}}
\toprule
Mathematical notation & Library notation \\
\midrule
$X$ connected & \texttt{isConnected 2 X} \\
$X$ simply connected & \texttt{isConnected 3 X} \\
$X$ $2$-connected & \texttt{isConnected 4 X} \\
$\tau_1X$ & \texttt{hLevelTrunc 3 X} \\
$\tau_2X$ & \texttt{hLevelTrunc 4 X} \\
\bottomrule
\end{tabular}
\end{center}
The source notation $\texttt{PT.}\lVert X\rVert_1$ denotes propositional
truncation. The declaration \texttt{Moore m k} counts suspensions from
$S^1$: it denotes $\Moore{m}{k+1}$ in this paper. In particular,
\texttt{Moore 2 1} and \texttt{Moore 10 2} denote $\Moore22$ and
$\Moore{10}3$.

\subsection{Statement-to-declaration correspondence}
Module paths below are relative to \path{agda/src/DeloopS3/}, except that
paths starting with \path{Cubical/} are relative to
\path{vendor/cubical/}. Declaration names include any enclosing named
parameterized module. Each listed file is in the active source tree.
\begingroup
\small
\setlength{\tabcolsep}{4pt}
\begin{longtable}{@{}>{\raggedright\arraybackslash}p{0.29\textwidth}>{\raggedright\arraybackslash}p{\dimexpr0.71\textwidth-8pt\relax}@{}}
\toprule
Paper statement & Exported declaration and module \\
\midrule
\endfirsthead
\toprule
Paper statement & Exported declaration and module (continued) \\
\midrule
\endhead
\bottomrule
\endfoot
Theorem~\ref{thm:main} & \texttt{Deloop-S3}\newline
  \path{Assembly/Complete.agda} \\[3pt]
Truncation square & \texttt{TruncationSquare.cartesian}\newline
  \path{Sphere/TruncationCartesian.agda} \\[3pt]
Hopf reflection & \texttt{hopfReflection}, \texttt{sphereUnitFiber}\newline
  \path{Sphere/HopfUnit.agda} \\[3pt]
Common fibre $F\simeq S^3$ & \texttt{RP\textsuperscript{2}-unit-fiber}\newline
  \path{Sphere/RP2UnitFiber.agda} \\[3pt]
Finite covers and total-space equivalences & \texttt{short-total}, \texttt{p}, \texttt{p-fibres}, \texttt{p-is-cover}\newline
  \path{Cert/Interpret/Done.agda} \\[3pt]
Twisted circle bundle & \texttt{Bundle.circleFibres}, \texttt{Bundle.baseFibre}\newline
  \path{Seifert/TwistedBundle/Cell.agda} \\[3pt]
Clutching coefficient $-2e$ in the calibrated chart & \texttt{Calibrated.coefficient}, \texttt{Calibrated.calibratedTotal}\newline
  \path{Seifert/TwistedBundle/TwoCap.agda} \\[3pt]
$\pi_2(M_e)=0$ for $e\ne0$ & \texttt{pi\textsubscript{2}-M\textsubscript{e}-zero}\newline
  \path{Seifert/TwistedBundle/Pi2Zero.agda} \\[3pt]
Dihedral-character rigidity & \texttt{dihedral-character-rigid}\newline
  \path{Bridge/SpanPi1/Character.agda} \\[3pt]
Source character & \texttt{source-character}\newline
  \path{Bridge/SpanPi1/SourceCharacter.agda} \\[3pt]
Nonabelian image and kernel witness & \texttt{subgroup-image-nonabelian}, \texttt{normal-core-lift}\newline
  \path{Bridge/SpanPi1/CoverWitnesses.agda} \\[3pt]
Identification of the kernel character & \texttt{cover-kernel-character-from-source}\newline
  \path{Bridge/SpanPi1/KernelCharacter.agda} \\[3pt]
Simple connectedness of the total pushouts & \texttt{X-simple}\newline
  \path{Bridge/SpanPi1/Target.agda} \\[3pt]
Cartesian span squares & \texttt{left-cartesian}, \texttt{right-cartesian}\newline
  \path{Assembly/MixedTruncation.agda} \\[3pt]
Fibre sequences $F\to E_c\to P_c$ & \texttt{local-fibre}\newline
  \path{Assembly/LocalCubes.agda} \\[3pt]
Actual cellular homology & \texttt{CW-AugChainComplex}, $\widetilde{\texttt H}^{\texttt{skel}}$\newline
  \path{Cubical/CW/ChainComplex.agda};\newline
  $\widetilde{\texttt H}^{\texttt{cw}}$\newline
  \path{Cubical/CW/Homology/Base.agda} \\[3pt]
Actual Hurewicz map and theorem & \texttt{HurewiczMapUntrunc}, \texttt{HurewiczTheorem}\newline
  \path{Cubical/CW/HurewiczTheorem.agda} \\[3pt]
Geometric exactness, $n<5$ & \texttt{SpanSequenceData}, \texttt{GeometricPushoutLES}\newline
  \path{Bridge/DMCExactness/Geometric.agda};\newline
  \texttt{finite-pushoutLES-merely}\newline
  \path{Bridge/CellularPushout/PrismChains/Packaging/Unconditional.agda} \\[3pt]
Boundary naturality, $H_3\to H_2$ & \texttt{TotalWindowBoundary.geometric-top-boundary-natural}\newline
  \path{Bridge/DegreeSpan/ContinuousBoundaryNatural.agda};\newline
  \texttt{initial-boundary-natural}\newline
  \path{Bridge/DegreeSpan/InitialBoundaryNatural.agda} \\[3pt]
Self-span naturality, $H_2\to H_1$ & \texttt{finite-same-span-boundary-natural}\newline
  \path{Bridge/DegreeSpan/LowerBoundaryDescent.agda} \\[3pt]
Component homology & \texttt{component-H\textsubscript{1}-groups},
  \texttt{source-H\textsubscript{2}-zero}, \texttt{target-H\textsubscript{2}-zero},
  \texttt{source-H\textsubscript{3}}, \texttt{target-H\textsubscript{3}}\newline
  \path{Bridge/Profile/ComponentHomologyComplete.agda} \\[3pt]
Covering-degree comparison & \texttt{comparisons}\newline
  \path{Bridge/DegreeSpan/LeftComparisons.agda};\newline
  \texttt{cover-degree-from-H2}\newline
  \path{Bridge/DegreeSpan/CoverDegreeFromH2.agda} \\[3pt]
Four-dimensional Moore recognition & \texttt{recognize-moore3}, \texttt{moore3-map-isEquiv}\newline
  \path{Bridge/Profile/MooreRecognition.agda} \\[3pt]
Bottom cofibre for the $I$ profile & \texttt{bottom-cofib-profile}\newline
  \path{Bridge/Profile/BottomCofibre.agda} \\[3pt]
Local contractions & \texttt{O-profile-contraction}, \texttt{I-profile-contraction}\newline
  \path{Bridge/Profile/Contraction.agda} \\[3pt]
Fibre preservation over simply connected bases & \texttt{localization-preserves-fiber}, \texttt{fiber-comparison}\newline
  \path{Local/RouteB2.agda} \\[3pt]
Fracture for $2$-connected types & \texttt{arithmetic-fracture-2-connected}\newline
  \path{Local/Fracture/Fracture.agda} \\[3pt]
Assembly of the delooping & \texttt{B-def}\newline
  \path{Assembly/DefB.agda};\newline
  \texttt{Deloop-final}\newline
  \path{Assembly/Final.agda} \\[3pt]
Recognition of the specified $3/10$ fracture cone & \texttt{unit-aware-fracture}\newline
  \path{HPInf/UnitFracture.agda} \\[3pt]
Cartesianness of \eqref{diag:HP} & \texttt{common-square-cartesian}\newline
  \path{HPInf/RelativeFracture.agda} \\[3pt]
Maximal torus (Theorem~\ref{thm:torus}) & \texttt{t}, \texttt{t-fiber}, \texttt{maximal-torus}\newline
  \path{HPInf/Complete.agda} \\
\end{longtable}
\endgroup

\subsection{Reproduction and development size}
\paragraph{Dependency identity.}
The official Cubical commit is\newline
\path{87af036ebd1c87956cf08fea44e62c7cf76392ca}, with parent\newline
\path{3d828afcac588414920565ad22d08d3dcdf02249}.
The full unmodified snapshot, including original licenses, is in
\path{vendor/cubical/}, locked by \path{agda/deps.lock.json}; official
provenance was verified through the GitHub API. The identifier
\texttt{cubical-0.9} does not mean release \texttt{v0.9}, which lacks the
required CW homology and Hurewicz modules. \path{agda/src/DeloopS3/Compat/Whitehead.agda}
reexports three declarations from \path{Cubical.Homotopy.Whitehead},
without copied mathematical types or proofs. Five files have import-only changes;
proof bodies and vendor sources are unchanged. The closure has 467 project and 416 official modules (883 total, excluding compiler builtins).

\paragraph{Reproduction.}
The checking scripts require Agda 2.8.0.2, Python 3.9 or later, Bash,
\texttt{flock}, and GNU \texttt{timeout}. From the root of a repository checkout, run:
\begin{quote}\small
\begin{verbatim}
bash agda/check.sh --configure-only
DELOOP_TIMEOUT=18000 bash agda/check.sh
\end{verbatim}
\end{quote}
The first command verifies the compiler version and dependency manifest
and creates a machine-local library registry. The second checks both
entries listed in \path{agda/roots.txt}. The runner uses this registry
and disables default libraries. It resolves paths relative to the
checkout, supports paths containing spaces, and performs no downloads.
Each check runs with one RTS capability and a 12\,GiB heap limit;
\texttt{DELOOP\_TIMEOUT} sets the per-entry time limit in seconds.

\paragraph{Verification.}
The official-dependency build began in a relocated directory containing
spaces with an empty user home and no project or library interfaces.
Five focused checks and both final entries exited with code zero;
successive stages reused only interfaces generated within that run.
The complete sequence took about 56 minutes, with 3,109 seconds for the
assembly stage and 54 seconds for the maximal-torus stage. Peak recorded
child-process RSS was approximately 12.4\,GiB under a 12\,GiB GHC heap cap.
The configuration and audit tools pass 48 regression tests.
The repository sources preserve all non-comment program lines of the checked
source; their only differences are comment and blank-line cleanup.
Both cleaned entry points subsequently passed an incremental
recheck with exit code zero.
Commands, hashes and measurements are retained in the author's local
verification record, \path{docs/verification/verification.json}; this file
is not included in the source repository.

\paragraph{Semantic audit.}
The author's local audit record, \path{docs/formalization-audit.md},
records the checked statement-to-source interfaces, assumptions, and coverage
limits separately from compilation and portability status. No target substitution was found in those
interfaces. This is not an exhaustive independent review of the deep
cubical coherence proofs.

\paragraph{Size.}
The active tree \path{agda/src} contains 467 Agda modules and 64,352 lines,
including comments and blank lines; \path{agda/archive} is excluded.
The two generated certificate-data modules
\path{Cert/Octa3.agda} and \path{Cert/Icosa10.agda} account for 453 of
these lines (27,105 bytes). The modules under \path{Cert/Input/} implement
the geometric interpretation of the input data and are included in the
remaining proof source. The totals can be reproduced from the repository
root by the following one-line command, whose output is the module count
followed by the line count.
\begin{quote}\scriptsize
\begin{verbatim}
find agda/src -name '*.agda' -exec wc -l {} + | awk '$2!="total"{n++;s+=$1}END{print n,s}'
\end{verbatim}
\end{quote}

\subsection{Finite certificates}
The certificate data are included as Agda declarations in
\path{agda/src/DeloopS3/Cert/Octa3.agda} and
\path{agda/src/DeloopS3/Cert/Icosa10.agda}. They specify permutation tables,
coset representatives, Schreier spanning trees, boundary words, filling
multiplicities, and the cancellation and shear data used in the short covers.
Declarations such as \texttt{group-ok}, \texttt{cosets-ok},
\texttt{graph-ok}, and \texttt{normalization-ok} validate the finite
computations by reduction against \path{Cert/Check.agda}.

The modules under \path{Cert/Input/} connect these checked data to the
covering diagrams. Those under \path{Cert/Interpret/} construct equivalences
of higher inductive types and the required boundary comparisons.
\path{Cert/Interpret/Done.agda} exports \texttt{short-total}, \texttt{p},
and \texttt{p-fibres}. Agda checks both the finite certificate assertions
and this semantic interpretation used by the final theorem.

\section*{Acknowledgements}
I thank Yan for providing the high-performance computing resources and the framework for collaboration among multiple large language models used in this project. I am especially grateful to Kang, who first introduced me to synthetic homotopy theory.

GPT and Claude were used extensively in the mathematical exploration, proof development, and formalization of this work. In particular, GPT suggested the strategy based on localization and arithmetic fracture that led to the construction presented here. Both GPT and Claude assisted with the formalization; Claude also assisted in coordinating this work and drafting the manuscript. I take full responsibility for the mathematical content, the formalization, and the exposition, including any errors or omissions.

\begin{thebibliography}{99}
\bibitem{hott} The Univalent Foundations Program. \emph{Homotopy Type Theory: Univalent Foundations of Mathematics}. 2013.
\bibitem{brunerie} G. Brunerie. \emph{On the homotopy groups of spheres in homotopy type theory}. PhD thesis, Universit\'e de Nice, 2016. arXiv:1606.05916.
\bibitem{cayley} U. Buchholtz and E. Rijke. The Cayley--Dickson construction in homotopy type theory. \emph{Higher Structures} 2(1), 2018, 30--41. arXiv:1610.01134.
\bibitem{highergroups} U. Buchholtz, F. van Doorn, and E. Rijke. Higher groups in homotopy type theory. \emph{LICS} 2018, 205--214. arXiv:1802.04315.
\bibitem{cors} J. D. Christensen, M. Opie, E. Rijke, and L. Scoccola. Localization in homotopy type theory. \emph{Higher Structures} 4(1), 2020, 1--32. arXiv:1807.04155.
\bibitem{scoccola} L. Scoccola. Nilpotent types and fracture squares in homotopy type theory. \emph{Math. Struct. Comput. Sci.} 30(5), 2020, 511--544. arXiv:1903.03245.
\bibitem{rector} D. L. Rector. Loop structures on the homotopy type of $S^3$. In \emph{Symposium on Algebraic Topology}, Lecture Notes in Math.\ 249, 1971, 99--105.
\bibitem{mcgibbon} C. A. McGibbon. Which group structures on $S^3$ have a maximal torus? In \emph{Geometric Applications of Homotopy Theory I}, Lecture Notes in Math.\ 657, 1978, 353--360.
\bibitem{dmw} W. G. Dwyer, H. R. Miller, and C. W. Wilkerson. The homotopic uniqueness of $BS^3$. In \emph{Algebraic Topology, Barcelona 1986}, Lecture Notes in Math.\ 1298, 1987, 90--105.
\bibitem{yaumaps} D. Yau. Maps to spaces in the genus of infinite quaternionic projective space. In \emph{Categorical Decomposition Techniques in Algebraic Topology}, Progr.\ Math.\ 215, Birkh\"auser, 2004, 293--302. arXiv:math/0209030.
\bibitem{yau} D. Yau. Loop structures on the homotopy type of $S^3$ revisited. \emph{Michigan Math.\ J.} 53, 2005, 283--290. arXiv:math/0501513.
\bibitem{cubicalagda} A. Vezzosi, A. M\"ortberg, and A. Abel. Cubical Agda: a dependently typed programming language with univalence and higher inductive types. \emph{Proc.\ ACM Program.\ Lang.} 3 (ICFP), article 87, 2019. doi:10.1145/3341691.
\bibitem{hurewiczhott} J. D. Christensen and L. Scoccola. The Hurewicz theorem in homotopy type theory. \emph{Algebr.\ Geom.\ Topol.} 23, 2023, 2107--2140. arXiv:2007.05833.
\bibitem{bm} K.-B. Hou (Favonia), E. Finster, D. R. Licata, and P. L. Lumsdaine. A mechanization of the Blakers--Massey connectivity theorem in homotopy type theory. \emph{LICS} 2016. arXiv:1605.03227.
\bibitem{cubical} The Cubical Agda Library. \url{https://github.com/agda/cubical}.
\end{thebibliography}
\end{document}
